\subsection{\texorpdfstring{Relations between \(\mathbf{S}(\mathbf{M}^{*})\), \(\mathbf{TS}(\mathbf{M})^{*\mathrm{gr}}\) and \(\operatorname{Pol}(\mathbf{M}, \mathbf{A})\)}{Relations between S(M*), the graded dual of TS(M), and Pol(M, A)}}

Let M be a free A-module, we shall equip the graded dual \(\mathbf{TS}(\mathbf{M})^{*gr}\) with the structure of a commutative, associative and unital graded \(^{1}\) algebra, derived from the graded cogebra structure of \(\mathbf{TS}(\mathbf{M})\) (III, p.~580). By III, p.~497 there exists a unique homomorphism of graded A-algebras

\[
\theta : \mathbf {S} (\mathbf {M} ^ {*}) \rightarrow \mathbf {T S} (\mathbf {M}) ^ {* \mathrm{gr}}
\]

inducing in degree 1 the identity mapping of \(\mathbf{M}^*\) .

PROPOSITION 20. --- If the A-module M is free and finitely generated, then \(\theta\) is an isomorphism of graded algebras.

Let \((e_i)_{i\in I}\) be a basis of \(\mathbf{M}\) and \((e_i^*)_{i\in I}\) the dual basis of \(\mathbf{M}^*\) . For \(\nu \in \mathbf{N}^I\) put

\[
e _ {\nu} = \prod_ {i \in I} \gamma_ {\nu_ {i}} (e _ {i}) \in \mathbf {T S} (\mathbf {M}).
\]

By Prop. 4 (IV, p.~47) the family \((e_{\nu})_{\nu \in \mathbf{N}^{I}}\) is a basis of \(\mathbf{TS}(\mathbf{M})\) ; let \((e_{\nu}^{*})\) be the basis of \(\mathbf{TS}(\mathbf{M})^{*\mathrm{gr}}\) dual to \((e_{\nu})\) . In the light of III, p.~505, Th. 1 it is enough to show that for any \(\nu \in \mathbf{N}^{I}\) we have

\[
e _ {\nu} ^ {*} = \prod_ {i \in I} \left(e _ {i} ^ {*}\right) ^ {\nu_ {i}},
\]

or also that for \(\rho, \sigma \in \mathbf{N}^{\mathrm{I}}\) we have \(e_{\rho}^{*} \cdot e_{\sigma}^{*} = e_{\rho + \sigma}^{*}\) ; but this last assertion follows from IV, p.~50, Formula (12).

Remark 1. --- In the same way we see that if M is a finitely generated free A-module, then the graded algebra \(\mathbf{S}(\mathbf{M})^{*\mathrm{gr}}\) defined in III, p.~593, may be identified with \(\mathbf{TS}(\mathbf{M}^*)\) .

PROPOSITION 21. --- The canonical homomorphism of A-modules (IV, p.~55)

\[
u: \mathbf {T S} (\mathrm{M}) ^ {* \mathrm{gr}} \rightarrow \operatorname{Pol} _ {\mathrm{A}} (\mathrm{M}, \mathrm{A})
\]

is an algebra homomorphism.

Let \(a \in TS^{q}(M)^{*}\) , \(b \in TS^{r}(M)^{*}\) , \(x \in M\) ; we have

\[
\begin{array}{r l} u (a b) (x) & = \langle a b, \gamma_ {q + r} (x) \rangle = \langle a \otimes b, c (\gamma_ {q + r} (x)) \rangle = \\ & = \langle a \otimes b, \gamma_ {q} (x) \otimes \gamma_ {r} (x) \rangle = \langle a, \gamma_ {q} (x) \rangle \langle b, \gamma_ {r} (x) \rangle = u (a) (x). u (b) (x), \end{array}
\]

whence the result.

Remarks. --- 2) The composite homomorphism \(\lambda_{\mathbf{M}}=u\circ\theta:\mathbf{S}(\mathbf{M}^{*})\to\mathrm{Pol}_{\mathbf{A}}(\mathbf{M},\mathbf{A})\) is the unique unital homomorphism of algebras inducing the inclusion of

\[
\mathbf {M} ^ {*} = \operatorname{Pol} ^ {1} (\mathbf {M}, \mathbf {A})
\]

in \(\operatorname{Pol}(\mathbf{M}, \mathbf{A})\) . If \(\mathbf{M}\) is finitely generated free and \(\mathbf{A}\) is an infinite integral domain, then \(\lambda_{\mathbf{M}}\) is bijective (Prop. 20 and Cor. of Prop. 19). In particular if \(\mathbf{A}\) is an infinite integral domain, then the canonical homomorphism \(f \mapsto \tilde{f}\) of \(\mathbf{A}[\mathbf{X}_1, \ldots, \mathbf{X}_n]\) into \(\operatorname{Pol}(\mathbf{A}^n, \mathbf{A})\) (IV, p.~4) is an isomorphism.

\begin{enumerate}
\def\labelenumi{\arabic{enumi})}
\setcounter{enumi}{2}
\tightlist
\item
  Consider the coproduct \(c_{\mathbf{S}}: \mathbf{S}(\mathbf{M}^{*}) \to \mathbf{S}(\mathbf{M}^{*} \times \mathbf{M}^{*})\) (III, p.~575, Example 6). For any \(v \in \mathbf{S}(\mathbf{M}^{*})\) , \(x, y \in \mathbf{M}\) , the polynomial mapping \(\lambda_{\mathbf{M} \times \mathbf{M}}(c_{\mathbf{S}}(v)): \mathbf{M} \times \mathbf{M} \to \mathbf{A}\) maps \((x, y)\) to \(\lambda_{\mathbf{M}}(v)(x + y)\) . For the two algebra homomorphisms of \(\mathbf{S}(\mathbf{M}^{*})\) into \(\operatorname{Map}(\mathbf{M} \times \mathbf{M}, \mathbf{A})\) defined in this way agree on \(\mathbf{M}^{*}\) , in virtue of the relation
\end{enumerate}

\[
(v \otimes 1 + 1 \otimes v) (x, y) = v (x + y) \quad (v \in \mathbf {M} ^ {*}).
\]

\section{SYMMETRIC FUNCTIONS}

\subsection{Symmetric polynomials}

Let \(n\) be a positive integer. For every permutation \(\sigma \in \mathfrak{S}_n\) let \(\varphi_{\sigma}\) be the automorphism of the A-algebra \(\mathbf{A}[\mathbf{X}_1, \ldots, \mathbf{X}_n]\) which maps \(\mathbf{X}_i\) to \(\mathbf{X}_{\sigma(i)}\) for \(1 \leqslant i \leqslant n\) . It is clear that \(\sigma \mapsto \varphi_{\sigma}\) is a homomorphism of \(\mathfrak{S}_n\) into the group of automorphisms of \(\mathbf{A}[\mathbf{X}_1, \ldots, \mathbf{X}_n]\) . We shall put \(\sigma f = \varphi_{\sigma}(f)\) for \(\sigma \in \mathfrak{S}_n\) and \(f \in \mathbf{A}[\mathbf{X}_1, \ldots, \mathbf{X}_n]\) . The polynomial \(f\) is said to be symmetric if \(\sigma f = f\) for all \(\sigma \in \mathfrak{S}_n\) ; the symmetric polynomials form a unital graded subalgebra of \(\mathbf{A}[\mathbf{X}_1, \ldots, \mathbf{X}_n]\) ; we shall denote it by \(\mathbf{A}[\mathbf{X}_1, \ldots, \mathbf{X}_n]^{\text{sym}}\) in the rest of this paragraph.

For every positive integer \(k\) we denote by \(\mathfrak{P}_k\) the set of \(k\) -element subsets of the set \(\{1, 2, \ldots, n\}\) and put

\[
s _ {k} = \sum_ {\mathrm{H} \in \mathfrak {P} _ {k}} \prod_ {i \in \mathrm{H}} X _ {i}.\tag{1}
\]

When we wish to specify the integer \(n\) we shall write \(s_{k,n}\) instead of \(s_k\) . We have in particular

\[
\begin{array}{l} s _ {0} = 1 \\ s _ {1} = \sum_ {1 \leqslant i \leqslant n} X _ {i} \\ s _ {2} = \sum_ {1 \leqslant i <   j \leqslant n} X _ {i} X _ {j} \\ \dots \dots \\ s _ {n} = X _ {1} \dots X _ {n} \end{array}
\]

and \(s_k = 0\) for \(k > n\) . It is clear that \(s_k\) is a homogeneous symmetric polynomial of degree \(k\) ; we shall call it the elementary symmetric polynomial of degree \(k\) .

In the ring \(\mathbf{A}[\mathbf{X}_1,\dots ,\mathbf{X}_n,\mathbf{U},\mathbf{V}]\) we have the relation

\[
\prod_ {i = 1} ^ {n} \left(\mathrm{U} + \mathrm{VX} _ {i}\right) = \sum_ {k = 0} ^ {n} \mathrm{U} ^ {n - k} \mathrm{V} ^ {k} s _ {k};\tag{2}
\]

by appropriate substitutions we deduce the relations

\[
\prod_ {i = 1} ^ {n} \left(1 + \mathrm{TX} _ {i}\right) = \sum_ {k = 0} ^ {n} s _ {k} \mathrm{T} ^ {k},\tag{3}
\]

\[
\prod_ {i = 1} ^ {n} \left(\mathbf {X} - \mathbf {X} _ {i}\right) = \sum_ {k = 0} ^ {n} (- 1) ^ {n - k} s _ {n - k} \mathbf {X} ^ {k}.\tag{4}
\]

THEOREM 1. --- Let us put \(E = A{[}X_{1}, \ldots, X_{n}{]}\) and \(S = A{[}X_{1}, \ldots, X_{n}{]}^{\text{sym}}\).

\begin{enumerate}
\def\labelenumi{\alph{enumi})}
\item
  The A-algebra S of symmetric polynomials is generated by \(s_1, \ldots, s_n\) .
\item
  The elements \(s_1, \ldots, s_n\) of \(\mathbf{E}\) are algebraically independent over \(\mathbf{A}\) (IV, p.~4).
\item
  The family of monomials \(\mathbf{X}^{\nu} = \mathbf{X}_1^{\nu(1)} \ldots \mathbf{X}_n^{\nu(n)}\) such that \(0 \leqslant \nu(i) < i\) for \(1 \leqslant i \leqslant n\) is a basis of the S-module \(\mathbf{E}\) . In particular, \(\mathbf{E}\) is a free S-module of rank \(n!\) .
\end{enumerate}

We shall prove the theorem by induction on \(n\) , the case \(n = 0\) being trivial. Let us write \(\mathbf{B} = \mathbf{A}[\mathbf{X}_n]\) and denote by \(s_k'\) the elementary symmetric polynomial of degree \(k\) in \(\mathbf{X}_1, \ldots, \mathbf{X}_{n-1}\) ; we thus have \(\mathbf{B}[\mathbf{X}_1, \ldots, \mathbf{X}_{n-1}] = \mathbf{A}[\mathbf{X}_1, \ldots, \mathbf{X}_n]\) . If we replace \(n\) by \(n-1\) and \(\mathbf{A}\) by \(\mathbf{B}\) in the statement of Th. 1, we can formulate the induction hypothesis as follows:

\begin{enumerate}
\def\labelenumi{(\Alph{enumi})}
\item
  The B-algebra \(S'\) of polynomials \(f \in A[X_1, \ldots, X_n]\) invariant under all permutations of \(X_1, \ldots, X_{n-1}\) is generated by \(s_1', \ldots, s_{n-1}'\) .
\item
  The elements \(s_1', \ldots, s_{n-1}'\) of \(E\) are algebraically independent over \(B\) .
\item
  The family of monomials \(\mathbf{X}_{1}^{\nu(1)}\ldots\mathbf{X}_{n-1}^{\nu(n-1)}\) such that \(0\leqslant\nu(i)<i\) for \(1\leqslant i\leqslant n-1\) is a basis of the S'-module E.
\end{enumerate}

We have the obvious relation

\[
s _ {k} = s _ {k} ^ {\prime} + s _ {k - 1} ^ {\prime} \mathbf {X} _ {n} \quad (1 \leqslant k \leqslant n - 1),\tag{5}
\]

from which we obtain by induction on \(k\)

\[
s _ {k} ^ {\prime} = (- 1) ^ {k} \mathbf {X} _ {n} ^ {k} + \sum_ {i = 1} ^ {k} (- 1) ^ {k - i} s _ {i} \mathbf {X} _ {n} ^ {k - i} \quad (1 \leqslant k \leqslant n - 1).\tag{6}
\]

We have \(\mathbf{S} \subset \mathbf{S}'\) , hence \(s_1, \ldots, s_n\) belong to \(\mathbf{S}'\) ; by (A) and formula (6) the B-algebra \(\mathbf{S}'\) is therefore generated by \(s_1, \ldots, s_{n-1}\) .

By (B) there exists an endomorphism \(u\) of the B-algebra \(S'\) such that

\[
u \left(s _ {k} ^ {\prime}\right) = (- 1) ^ {k} X _ {n} ^ {k} + \sum_ {i = 1} ^ {k} (- 1) ^ {k - i} s _ {i} ^ {\prime} X _ {n} ^ {k - i} \quad (1 \leqslant k \leqslant n - 1).\tag{7}
\]

By (5), we have \(u(s_k) = u(s_k') + u(s_{k-1}') \mathbf{X}_n\) , whence \(u(s_k) = s_k'\) by an easy calculation. Let \(\mathbf{P} \in \mathbf{B}[\mathbf{Y}_1, \ldots, \mathbf{Y}_{n-1}]\) ; then from \(\mathbf{P}(s_1, \ldots, s_{n-1}) = 0\) we deduce

\[
0 = u \left(\mathrm{P} \left(s _ {1}, \dots , s _ {n - 1}\right)\right) = \mathrm{P} \left(s _ {1} ^ {\prime}, \dots , s _ {n - 1} ^ {\prime}\right),
\]

whence P = 0 by (B). It follows that the B-algebra \(S'\) is generated by the algebraically independent elements \(s_{1}, \ldots, s_{n-1}\) . This property may be reformulated as follows:

\begin{enumerate}
\def\labelenumi{(\Alph{enumi})}
\setcounter{enumi}{3}
\tightlist
\item
  The A-algebra \(S'\) is generated by the algebraically independent elements \(s_1, \ldots, s_{n-1}, X_n\) .
\end{enumerate}

We can thus identify \(S'\) with the polynomial ring \(C[X_n]\) , where \(C\) is the sub-A-algebra of \(E\) generated by \(s_1, \ldots, s_{n-1}\) .

To prove a), let \(f \in S\) be a homogeneous symmetric polynomial of degree \(m\) . We have \(f \in S' = C[X_n]\) , so there exists an element \(g = P(s_1, \ldots, s_{n-1})\) of \(C\) , homogeneous of degree \(m\) in \(X_1, \ldots, X_n\) , such that \(f - g\) is divisible by \(X_n\) . Since \(f - g\) is symmetric, each of its terms is also divisible by \(X_1, \ldots, X_{n-1}\) , hence \(f - g\) is divisible by \(s_n = X_1 \ldots X_n\) . In other words, there exists \(h \in S\) such that \(f = g + hs_n\) , whence \(\deg h < m\) . By induction on \(m\) it follows that \(f\) belongs to

\[
\mathbf {C} \left[ s _ {n} \right] _ {\mathrm{E}} = \mathbf {A} \left[ s _ {1}, \dots , s _ {n - 1}, s _ {n} \right] _ {\mathrm{E}}.
\]

Thus we have shown that the A-algebra S is generated by \(s_{1}, \ldots, s_{n}\) .

Next we prove \(b\) ). If we substitute \(\mathbf{X}_n\) for \(\mathbf{X}\) in (4), we find

\[
(- 1) ^ {n + 1} s _ {n} = \mathrm{X} _ {n} ^ {n} + \sum_ {k = 1} ^ {n - 1} (- 1) ^ {n - k} s _ {n - k} \mathrm{X} _ {n} ^ {k};
\]

in other words, \((-1)^{n + 1}s_n\) is a monic polynomial in \(X_{n}\) of degree \(n\) and with coefficients in C. By IV, p.~11 we thus have the following property:

\begin{enumerate}
\def\labelenumi{(\Alph{enumi})}
\setcounter{enumi}{4}
\tightlist
\item
  The C-algebra homomorphism \(\varphi\) of C{[}T{]} into C{[}X \(_n\) {]} = S' which maps T to \(s_n\) is injective, and S' is a free module over the image of \(\varphi\) , with basis (1, X \(_n\) , \ldots, X \(_n^{n-1}\) ).
\end{enumerate}

The elements \(s_{1}, \ldots, s_{n-1}\) of C are algebraically independent over A, by (D); thus the injectivity of \(\varphi\) means that \(s_{1}, \ldots, s_{n-1}, s_{n}\) are algebraically independent over A, whence b).

To prove c), the image of \(\varphi\) equals \(C[s_{n}]_{E} = S\) , hence by (E), \(S'\) is a free module over S on the basis \((1, X_{n}, ..., X_{n}^{n-1})\) . Now the assertion c) follows from the induction hypothesis (C) and Prop. 25 of II, p.~222. This completes the proof.

Let \(f\) be a symmetric polynomial in \(\mathbf{X}_1, \ldots, \mathbf{X}_n\) , homogeneous of degree \(m\) . By Th. 1 (IV, p.~62) there exists a polynomial \(\mathbf{Q} \in \mathbf{A}[\mathbf{Y}_1, \ldots, \mathbf{Y}_n]\) such that \(f = \mathbf{Q}(s_1, \ldots, s_n)\) . The preceding proof provides a means of explicit calculation of \(\mathbf{Q}\) , by a double induction on \(n\) and \(m\) . For as we have seen, there exists a polynomial \(\mathbf{P} \in \mathbf{A}[\mathbf{Y}_1, \ldots, \mathbf{Y}_{n-1}]\) and a polynomial \(h\) symmetric in \(\mathbf{X}_1, \ldots, \mathbf{X}_n\) which is homogeneous of degree \(m - n\) , such that

\[
f = \mathrm{P} (s _ {1}, \dots , s _ {n - 1}) + s _ {n} h.\tag{8}
\]

For every polynomial \(u \in \mathbf{A}[\mathbf{X}_1, \ldots, \mathbf{X}_n]\) we put

\[
u ^ {\prime} \left(\mathrm{X} _ {1}, \dots , \mathrm{X} _ {n - 1}\right) = u \left(\mathrm{X} _ {1}, \dots , \mathrm{X} _ {n - 1}, 0\right).
\]

Then \(s_1', \ldots, s_{n-1}'\) are the elementary symmetric polynomials in \(X_1, \ldots, X_{n-1}\) and formula (8) implies

\[
f ^ {\prime} = \mathrm{P} \left(s _ {1} ^ {\prime}, \dots , s _ {n - 1} ^ {\prime}\right).
\]

The determination of P is thus reduced to a calculation of symmetric polynomials in n - 1 indeterminates, and h is obtained from (8).

We illustrate the method by two examples.

Examples. --- 1) Let \(n = 3\) and

\[
f = \mathbf {X} _ {1} ^ {2} (\mathbf {X} _ {2} + \mathbf {X} _ {3}) + \mathbf {X} _ {2} ^ {2} (\mathbf {X} _ {3} + \mathbf {X} _ {1}) + \mathbf {X} _ {3} ^ {2} (\mathbf {X} _ {1} + \mathbf {X} _ {2}).
\]

We have

\[
f ^ {\prime} = \mathbf {X} _ {1} ^ {2} \mathbf {X} _ {2} + \mathbf {X} _ {1} \mathbf {X} _ {2} ^ {2} = \mathbf {X} _ {1} \mathbf {X} _ {2} (\mathbf {X} _ {1} + \mathbf {X} _ {2}) = s _ {1} ^ {\prime} s _ {2} ^ {\prime}.
\]

Let us write \(g = f - s_1s_2\) ; we have

\[
\boldsymbol {g} = f - \left(\mathbf {X} _ {1} + \mathbf {X} _ {2} + \mathbf {X} _ {3}\right) \left(\mathbf {X} _ {1} \mathbf {X} _ {2} + \mathbf {X} _ {1} \mathbf {X} _ {3} + \mathbf {X} _ {2} \mathbf {X} _ {3}\right) = - 3 \mathbf {X} _ {1} \mathbf {X} _ {2} \mathbf {X} _ {3},
\]

hence finally

\[
f = s _ {1} s _ {2} - 3 s _ {3}.
\]

\begin{enumerate}
\def\labelenumi{\arabic{enumi})}
\setcounter{enumi}{1}
\tightlist
\item
  Again let \(n = 3\) and put \(p = \mathbf{X}_1^3 + \mathbf{X}_2^3 + \mathbf{X}_3^3\) .
\end{enumerate}

We have \(p(\mathbf{X}_1, 0, 0) = \mathbf{X}_1^3 = s_1(\mathbf{X}_1, 0, 0)^3\) . Writing \(q = p - s_1^3\) , we obtain

\[
q = - 3 f - 6 \mathrm{X} _ {1} \mathrm{X} _ {2} \mathrm{X} _ {3} = - 3 s _ {1} s _ {2} + 3 s _ {3}
\]

and finally

\[
p = s _ {1} ^ {3} - 3 s _ {1} s _ {2} + 3 s _ {3}.
\]

Let \(S_1, \ldots, S_n\) be indeterminates. We shall equip the polynomial algebra \(A[S_1, \ldots, S_n]\) with the graduation of type N for which \(S_k\) is of weight k for \(1 \leqslant k \leqslant n\) (IV, p.~3), and we equip \(A[X_1, \ldots, X_n]\) with the usual graduation. For \(1 \leqslant k \leqslant n\) the elementary symmetric polynomial \(s_{k,n}\) in \(X_1, \ldots, X_n\) is homogeneous of degree k. By Th. 1 (IV, p.~62) the mapping \(g \mapsto g(s_{1,n}, \ldots, s_{n,n})\) is thus an isomorphism of graded algebras

\[
\varphi_ {n}: \mathrm{A} [ \mathrm{S} _ {1}, \dots , \mathrm{S} _ {n} ] \rightarrow \mathrm{A} [ \mathrm{X} _ {1}, \dots , \mathrm{X} _ {n} ] ^ {\text { sym }}.
\]

Let \(m\) be an integer such that \(0 \leqslant m \leqslant n\) . For every integer \(k\) such that \(1 \leqslant k \leqslant m\) we have

\[
s _ {k, m} (\mathbf {X} _ {1}, \dots , \mathbf {X} _ {m}) = s _ {k, n} (\mathbf {X} _ {1}, \dots , \mathbf {X} _ {m}, 0, \dots , 0)\tag{9}
\]

by definition (IV, p.~61, formula (1)) of \(s_k\) . Hence the diagram

\[
\begin{array}{c} \mathrm{A} [ \mathrm{S} _ {1}, \dots , \mathrm{S} _ {m} ] \xrightarrow {j} \mathrm{A} [ \mathrm{S} _ {1}, \dots , \mathrm{S} _ {n} ] \\ \varphi_ {m} \Bigg | _ {\downarrow} \qquad \qquad \qquad \qquad \qquad \qquad \qquad \qquad \qquad \qquad \varphi_ {n} \\ \mathrm{A} [ \mathrm{X} _ {1}, \dots , \mathrm{X} _ {m} ] ^ {\text { sym }} \xleftarrow {\rho} \mathrm{A} [ \mathrm{X} _ {1}, \dots , \mathrm{X} _ {n} ] ^ {\text { sym }} \end{array}\tag{10}
\]

(where \(j\) denotes the canonical inclusion and \(\rho\) the homomorphism

\[
g \mapsto g \left(\mathrm{X} _ {1}, \dots , \mathrm{X} _ {m}, 0, \dots , 0\right))
\]

is commutative.

PROPOSITION 1. --- For every pair of positive integers \(k\) , \(n\) , let \(\mathbf{S}_{k}^{(n)}\) be the A-module consisting of all symmetric polynomials in \(\mathbf{X}_{1},\ldots,\mathbf{X}_{n}\) , which are homogeneous of degree \(k\) . If the integer \(m\) satisfies \(0\leqslant k\leqslant m\leqslant n\) , then the mapping \(f\mapsto f(\mathbf{X}_{1},\ldots,\mathbf{X}_{m},0,\ldots,0)\) is an isomorphism of \(\mathbf{S}_{k}^{(n)}\) onto \(\mathbf{S}_{k}^{(m)}\) .

By the commutativity of the diagram (10) it is enough to show that every isobaric polynomial of weight k in \(S_{1}, \ldots, S_{n}\) depends only on \(S_{1}, \ldots, S_{m}\) under the hypotheses \(0 \leqslant k \leqslant m \leqslant n\) . Now the weight of a monomial \(\mathrm{S}_{1}^{\alpha(1)} \ldots \mathrm{S}_{n}^{\alpha(n)}\) is equal to the integer \(\alpha(1) + 2\alpha(2) + \cdots + n\alpha(n)\) ; since the integers \(\alpha(1), \ldots, \alpha(n)\) are positive, the relation

\[
\alpha (1) + 2 \alpha (2) + \dots + n \alpha (n) = k \leqslant n
\]

implies \(\alpha(j) = 0\) for \(k < j \leqslant n\) , whence the assertion.

Example 3. --- By example 2 of IV, p.64 and Prop. 1 above we thus have

\[
\sum_ {i = 1} ^ {n} X _ {i} ^ {3} = s _ {1, n} ^ {3} - 3 s _ {1, n} s _ {2, n} + 3 s _ {3, n}
\]

for every integer \(n \geqslant 3\) . The commutativity of the diagram (10) moreover gives the formulae

\[
\begin{array}{c} \mathbf {X} _ {1} ^ {3} + \mathbf {X} _ {2} ^ {3} = s _ {1, 2} ^ {3} - 3 s _ {1, 2} s _ {2, 2}, \\ \mathbf {X} _ {1} ^ {3} = s _ {1, 1} ^ {3}. \end{array}
\]

Remark. --- Let n and k be two positive integers. We denote by \(\Delta_{k,n}\) the set of elements of length k in \(N^{n}\) and further equip \(\Delta_{k,n}\) with the order relation, written \(\alpha \preccurlyeq \beta\) , induced by the lexicographic order on \(N^{n}\) (Set Theory, III, p.~157), and define an action of the group \(S_{n}\) on \(N^{n}\) by \((\sigma\alpha)(i) = \alpha(\sigma^{-1}(i))\) for \(\sigma \in S_{n}\) , \(\alpha \in N^{n}\) and \(1 \leqslant i \leqslant n\) . Moreover, we denote by \(D_{k}\) the set of elements \(\alpha = (\alpha(1), \ldots, \alpha(k))\) of \(N^{k}\) such that

\[
\alpha (1) \geqslant \alpha (2) \geqslant \dots \geqslant \alpha (k), \quad \alpha (1) + \dots + \alpha (k) = k.
\]

Suppose that \(k \leqslant n\) and let us identify \(\mathbf{N}^k\) with a subset of \(\mathbf{N}^n\) by the mapping \((\alpha(1), \ldots, \alpha(k)) \mapsto (\alpha(1), \ldots, \alpha(k), 0, \ldots, 0)\) . Then \(D_k\) consists of the elements \(\alpha\) of \(\Delta_{k,n}\) such that \(\sigma \alpha \preccurlyeq \alpha\) for all \(\sigma \in \mathfrak{S}_n\) . It follows that every orbit of the group \(\mathfrak{S}_n\) in \(\Delta_{k,n}\) contains a unique element of \(D_k\) . For every \(\alpha \in D_k\) let \(O(\alpha)\) be the orbit of \(\alpha\) in \(\Delta_{k,n}\) under the operation of \(\mathfrak{S}_n\) ; put

\[
\mathbf {M} (\alpha) = \sum_ {\beta \in O (\alpha)} \mathbf {X} ^ {\beta}.\tag{11}
\]

It follows from Lemma 1 of IV, p.~47, that the family \((\mathbf{M}(\alpha))_{\alpha \in \mathrm{D}_k}\) is a basis of the A-module \(\mathbf{S}_k^{(n)}\) .

For each \(\alpha \in D_k\) put

\[
\mathrm{S} (\alpha) = \prod_ {i = 1} ^ {k} s _ {i} ^ {\alpha (i) - \alpha (i + 1)} \quad (\text { where   it   is   understood   that } \alpha (k + 1) = 0);\tag{12}
\]

since we have \(\sum_{i=1}^{k} i \cdot (\alpha(i) - \alpha(i+1)) = \sum_{i=1}^{k} \alpha(i) = k\) , the symmetric polynomial \(S(\alpha)\) is homogeneous of degree \(k\) . It follows directly from Th. 1 (IV, p.~62) that the family \((S(\alpha))_{\alpha \in D_k}\) is a basis of the A-module \(S_k^{(n)}\) .

Let \(\alpha, \beta\) be in \(D_k\) , and let \(c_{\alpha\beta}\) be the coefficient of the monomial \(X^\beta\) in the polynomial \(S(\alpha)\) defined by (12) in the case \(A = Z\) and \(k = n\) . This is then a positive integer, independent of the ring \(A\) and of the integer \(n\) . By formula (9) (IV, p.~64) we then have

\[
\mathrm{S} (\alpha) = \sum_ {\beta \in \mathrm{D} _ {k}} c _ {\alpha \beta} \cdot \mathrm{M} (\beta) \quad (\alpha \in \mathrm{D} _ {k}).\tag{13}
\]

It may be shown (cf.~IV, p.~101, Exercise 13) that the matrix \(\mathbf{C} = (c_{\alpha \beta})_{\alpha, \beta \in D_k}\) is such that \(c_{\alpha \alpha} = 1\) and \(c_{\alpha \beta} = 0\) for \(\alpha < \beta\) ; generalizing the terminology introduced in II, p.~351, we shall say that \(\mathbf{C}\) belongs to the lower total triangular group. The same is true of the matrix \(\mathbf{D}\) inverse to \(\mathbf{C}\) . The values of the matrices \(\mathbf{C}\) and \(\mathbf{D}\) when \(2 \leqslant k \leqslant 5\) , are given in the table (IV, p.~103-105).

Suppose now that \(n < k\) and let us identify \(\mathbf{N}^n\) with a subset of \(\mathbf{N}^k\) by the mapping \((\alpha(1), \ldots, \alpha(n)) \mapsto (\alpha(1), \ldots, \alpha(n), 0, \ldots, 0)\) . For every \(\alpha\) in \(D_k \cap N^n\) we again denote by \(O(\alpha)\) the orbit of \(\alpha\) in \(\Delta_{k,n}\) under the operation of \(\mathfrak{S}_n\) and define \(M(\alpha)\) by (11). Further we define \(S(\alpha)\) by (12), then the families \((M(\alpha))_{\alpha \in D_k \cap N^n}\) and \((S(\alpha))_{\alpha \in D_k \cap N^n}\) are bases of the A-module \(S_k^{(n)}\) . By formula (9) of IV, p.~64, we have a formula analogous to (13) where \(D_k\) is replaced by \(D_k \cap N^n\) , with the same integers \(c_{\alpha\beta}\) .

Example 4. --- By Example 3 of IV, p.~65, we have

\[
\mathbf {M} (3, 0, 0) = \mathbf {S} (3, 0, 0) - 3 \mathbf {S} (2, 1, 0) + 3 \mathbf {S} (1, 1, 1)
\]

for every integer \(n \geqslant 3\) , and so

\[
\mathbf {M} (3, 0) = \mathbf {S} (3, 0) - 3 \mathbf {S} (2, 1)
\]

for n = 2.

\subsection{Symmetric rational fractions}

Let K be a commutative field and \(X_{1}, X_{2}, \ldots, X_{n}\) indeterminates. For every \(\sigma \in S_{n}\) we have defined in No.~1 (IV, p.~61) an automorphism \(\varphi_{\sigma}\) of \(K[X_{1}, X_{2}, \ldots, X_{n}]\) . This automorphism has a unique extension \(\psi_{\sigma}\) to the field \(K(X_{1}, \ldots, X_{n})\) , and \(\sigma \mapsto \psi_{\sigma}\) is an injective homomorphism of \(S_{n}\) into the automorphism group of \(K(X_{1}, \ldots, X_{n})\) . For each \(f \in K(X_{1}, \ldots, X_{n})\) we have \((\psi_{\sigma} f) (X_{1}, \ldots, X_{n}) = f(X_{\sigma(1)}, \ldots, X_{\sigma(n)})\) . The rational fractions f such that \(\psi_{\sigma}(f) = f\) for all \(\sigma \in S_{n}\) are called symmetric rational fractions. The set of all symmetric rational fractions in \(X_{1}, \ldots, X_{n}\) is a subfield of \(K(X_{1}, \ldots, X_{n})\) .

PROPOSITION 2. --- The field of symmetric rational fractions in \(X_{1}, \ldots, X_{n}\) is the field of fractions of the ring of symmetric polynomials in \(X_{1}, \ldots, X_{n}\) .

Let \(f \in \mathbf{K}(\mathbf{X}_1, \ldots, \mathbf{X}_n)\) be a symmetric rational fraction, and let \(u_1, v_1\) be two elements of \(\mathbf{K}[\mathbf{X}_1, \ldots, \mathbf{X}_n]\) such that \(f = \frac{u_1}{v_1}\) . We put \(v = \prod_{\sigma \in \mathfrak{S}_n} \psi_\sigma(v_1) \in \mathbf{K}[\mathbf{X}_1, \ldots, \mathbf{X}_n]\) and \(u = fv \in \mathbf{K}[\mathbf{X}_1, \ldots, \mathbf{X}_n]\) ; then \(v\) is symmetric, hence \(u\) is symmetric, because \(f\) is, and \(f = \frac{u}{v}\) , whence the result.

COROLLARY. --- Let \(s_{1}, s_{2}, \ldots, s_{n}\) be the elementary symmetric polynomials in \(X_{1}, \ldots, X_{n}\) . For every rational fraction \(g \in \mathbf{K}(S_{1}, S_{2}, \ldots, S_{n})\) the sequence \((s_{1}, s_{2}, \ldots, s_{n})\) is substitutable in g and the mapping \(g \mapsto g(s_{1}, s_{2}, \ldots, s_{n})\) is an isomorphism of \(\mathbf{K}(S_{1}, S_{2}, \ldots, S_{n})\) onto the field of symmetric rational fractions in \(X_{1}, \ldots, X_{n}\) .

This follows from Prop. 2 and Th. 1 of IV, p.~62.

\subsection{Symmetric formal power series}

Let I be a set, \(\mathbf{X} = (\mathbf{X}_i)_{i\in I}\) a family of indeterminates and \(\mathbf{A}[[\mathbf{X}]]\) the algebra of formal power series in the \(\mathbf{X}_i\) . For every permutation \(\sigma \in \mathfrak{S}_I\) there exists a unique continuous automorphism \(\varphi_{\sigma}\) of the algebra \(\mathbf{A}[[\mathbf{X}]]\) which maps \(\mathbf{X}_i\) to \(\mathbf{X}_{\sigma (i)}\) for each \(i\in I\) (IV, p.~28, Prop. 4); it is clear that \(\sigma \mapsto \varphi_{\sigma}\) is a homomorphism of \(\mathfrak{S}_I\) into the group of continuous automorphisms of the algebra \(\mathbf{A}[[\mathbf{X}]]\) . Let \(f\in \mathbf{A}[[\mathbf{X}]]\) be a formal power series; we put \(\sigma f = \varphi_{\sigma}(f)\) and we shall say that the formal power series \(f\) is symmetric if \(\sigma f = f\) for each \(\sigma \in \mathfrak{S}_I\) . The symmetric formal power series form a closed subalgebra of \(\mathbf{A}[[\mathbf{X}]]\) which is denoted by \(\mathbf{A}[[\mathbf{X}]]^{\mathrm{sym}}\) and is equipped with the topology induced by that of \(\mathbf{A}[[\mathbf{X}]]\) .

Let T be an indeterminate. In the ring of formal power series A{[}{[}X, T{]}{]} the family \((\mathbf{X}_i\mathbf{T})_{i\in I}\) is summable, hence the family \((1 + \mathbf{X}_i\mathbf{T})_{i\in I}\) is multipliable (IV, p.~27, Prop. 2); moreover we have

\[
\prod_ {i \in I} (1 + X _ {i} T) = 1 + \sum_ {k \geqslant 1} s _ {k} T ^ {k},\tag{14}
\]

where the formal power series \(s_k \in \mathbf{A}[[\mathbf{X}]]\) is defined by

\[
s _ {k} = \sum_ {\mathrm{H} \in \mathfrak {P} _ {k}} \left(\prod_ {i \in \mathrm{H}} \mathrm{X} _ {i}\right) \quad (k \geqslant 1)\tag{15}
\]

(we denote by \(\mathfrak{P}_k\) the set of all \(k\) -element subsets of I). In particular we have \(s_1 = \sum_{i \in I} X_i\) . When I is finite with \(n\) elements we have \(s_k = 0\) for \(k > n\) ; more precisely, when \(\mathbf{I} = \{1,\dots,n\}\) then the formal power series \(s_k\) is nothing other than the elementary symmetric polynomial of degree \(k\) in \(\mathbf{X}_1,\ldots,\mathbf{X}_n\) .

Let \(\mathbf{S} = (\mathbf{S}_k)_{k\geqslant 1}\) be a sequence of indeterminates. Since the formal power series \(s_k\) is of order \(\geqslant k\) , and belongs to \(\mathbf{A}[[\mathbf{X}]]^{\mathrm{sym}}\) , conditions a) and b) of Prop. 4 of IV, p.~28 are satisfied with \(\mathrm{E} = \mathbf{A}[[\mathbf{X}]]^{\mathrm{sym}}\) ; there exists thus a unique continuous A-algebra homomorphism

\[
\varphi_ {I}: \mathrm{A} [ [ \mathrm{S} ] ] \rightarrow \mathrm{A} [ [ \mathrm{X} ] ] ^ {\text { sym }}
\]

such that \(\varphi_{I}(S_{k}) = s_{k}\) for each integer \(k\geqslant 1\)

THEOREM 2. --- a) If I is a finite set of n elements, then \(\varphi_{I}\) induces a bicontinuous isomorphism of A{[}{[}S \(_{1}\) , \ldots, S \(_{n}\) {]}{]} onto A{[}{[}X{]}{]} \(^{\text{sym}}\) .

\begin{enumerate}
\def\labelenumi{\alph{enumi})}
\setcounter{enumi}{1}
\tightlist
\item
  If I is infinite, \(\varphi_{\mathrm{I}}\) is a bicontinuous isomorphism of A{[}{[}S{]}{]} onto A{[}{[}X{]}{]} \(^{\mathrm{sym}}\) .
\end{enumerate}

In case a) we put \(\mathbf{B} = \mathbf{A}[[\mathbf{S}_1, \ldots, \mathbf{S}_n]]\) and equip this algebra with the topology induced by that of \(\mathbf{A}[[\mathbf{S}]]\) ; we also denote by \(\psi_{\mathrm{I}}\) the restriction of \(\varphi_{\mathrm{I}}\) to \(\mathbf{B}\) . In case b) we put \(\mathbf{B} = \mathbf{A}[[\mathbf{S}]]\) and \(\psi_{\mathrm{I}} = \varphi_{\mathrm{I}}\) . We shall equip the polynomial algebra \(\mathbf{A}[\mathbf{S}]\) with the graduation of type \(\mathbf{N}\) for which \(\mathbf{S}_k\) is of weight \(k\) for every integer \(k \geqslant 1\) .

Lemma 1. --- Let J be a finite subset of I, r an integer such that Card \(J \geqslant r\) and f a symmetric formal power series which is homogeneous of degree r in the \(X_{i}\) ( \(i \in I\) ). Let \(\bar{f}\) be the formal power series obtained by substituting 0 for \(X_{i}\) for each i in I - J. If \(\bar{f} = 0\) , then we have f = 0.

Let us put \(f = \sum_{|\alpha| = r} a_{\alpha} X^{\alpha}\) (where \(|\alpha|\) is the length \(\sum_{i \in I} \alpha_i\) of the multiindex \(\alpha = (\alpha_i)_{i \in I}\) ). If \(\alpha\) is a multiindex of length \(r\) , and \(J'\) is the support of \(\alpha\) (the set of \(i \in I\) such that \(\alpha_i \neq 0\) ), then Card \(J' \leqslant r\) , hence there exists a permutation \(\sigma \in \mathfrak{S}_I\) such that \(\sigma(J') \subset J\) . Put \(\beta_i = \alpha_{\sigma^{-1}(i)}\) for \(i \in I\) , then the monomial \(X^\beta = \prod_{i \in I} X_{\sigma(i)}^{\alpha_i}\) depends only on the indeterminates \(X_j (j \in J)\) , whence \(a_\beta = 0\) by the hypothesis \(\bar{f}=0\) . Since f is symmetric, we have \(a_{\alpha}=a_{\beta}\) and since \(\alpha\) was arbitrary, we conclude that f=0.

Lemma 2. --- Let \(f\) be a symmetric formal power series of degree \(r\) in the \(X_i (i \in I)\) . There exists a unique polynomial \(P \in B \cap A[S]\) , isobaric of weight \(r\) , such that \(f = \psi_I(P)\) .

The case when I is finite follows from Th. 1 (IV, p.~62).

Assume that I is infinite and choose a finite subset J of I, containing r elements. We keep the notation of Lemma 1; we remark that every isobaric polynomial of weight r in the \(S_{n}\) ( \(n \geqslant 1\) ) depends only on \(S_{1}, \ldots, S_{r}\) and that \(\bar{s}_{1}, \ldots, \bar{s}_{r}\) are the elementary symmetric polynomials in the r indeterminates \(X_{j}\) ( \(j \in J\) ). If P is an isobaric polynomial of weight r in the \(S_{n}\) and \(h = f - \psi_{\mathrm{I}}(\mathrm{P})\) , then \(\bar{h} = \bar{f} - \mathrm{P}(\bar{s}_{1}, \ldots, \bar{s}_{r})\) and Lemma 1 shows the relation \(f = \psi_{\mathrm{I}}(\mathrm{P})\) to be equivalent to \(\bar{f} = \mathrm{P}(\bar{s}_{1}, \ldots, \bar{s}_{r})\) . By Th. 1 (IV, p.~62) there exists a unique polynomial \(P \in A[S]\) , isobaric of weight r and such that \(\bar{f} = \mathrm{P}(\bar{s}_{1}, \ldots, \bar{s}_{r})\) , whence the result.

Lemma 3. --- For every integer \(m \geqslant 0\) let \(\mathfrak{c}_m\) be the ideal of the algebra \(\mathrm{A}[[\mathrm{X}]]^{\mathrm{sym}}\) consisting of all symmetric formal power series of order \(\geqslant m\) . The sequence \((\mathfrak{c}_m)_{m \geqslant 0}\) is a fundamental system of neighbourhoods of 0 in \(\mathrm{A}[[\mathrm{X}]]^{\mathrm{sym}}\) .

The lemma clearly holds when I is finite, so let us assume that I is infinite. For every finite subset J of I, consisting of m elements, denote by \(\tilde{J}\) the set of elements of \(\mathbf{N}^{(1)}\) of length \textless{} m and with support in J. Further denote by \(a_{j}^{\prime}\) the set of formal power series containing no term \(aX^{\alpha}\) with \(\alpha \in \tilde{J}\) . Since \(\tilde{J}\) is a finite subset of \(\mathbf{N}^{(1)}\) and every finite subset of \(\mathbf{N}^{(1)}\) is contained in a set of the form \(\tilde{J}\) , the family \((\mathfrak{a}_{\mathrm{j}}^{\prime})\) is a base of neighbourhoods of 0 in A{[}{[}X{]}{]} (IV, p.~26). Now Lemma 1 implies the relation \(a_{j}^{\prime} \cap A[[X]]^{\text{sym}} = c_{m}\) for every subset J with m elements, and this proves Lemma 3.

Since there are only a finite number of monomials of a given weight in the \(S_{k}\) , every formal power series \(f \in B\) may be written in a unique way as \(f = \sum_{r \geqslant 0} P_{r}\) , where \(P_{r}\) is an isobaric polynomial of weight r in \(B \cap A[S]\) . For every integer \(m \geqslant 0\) let \(\mathfrak{b}_m\) be the ideal of B consisting of all formal power series of the above type such that \(P_r = 0\) for \(0 \leqslant r < m\) . The sequence \((\mathfrak{b}_m)_{m \geqslant 0}\) is a base of neighbourhoods of 0 in B.

With the above notation \(\psi_{\mathrm{I}}(\mathbf{P}_{r})\) is a symmetric formal power series in the \(\mathbf{X}_i\) , homogeneous of degree \(r\) , and so this is the homogeneous component of degree \(r\) of \(\psi_{\mathrm{I}}(f)\) . Lemma 2 shows that \(\psi_{\mathrm{I}}\) is an algebra homomorphism of B onto \(\mathbf{A}[[\mathbf{X}]]^{\mathrm{sym}}\) , mapping \(\mathfrak{b}_m\) to \(\mathfrak{c}_m\) for every integer \(m \geqslant 0\) ; now Lemma 3 shows \(\psi_{\mathrm{I}}\) to be bicontinuous, and this completes the proof of Theorem 2.

\subsection{Sums of powers}

We again write \(\mathbf{X} = (\mathbf{X}_{i})_{i \in I}\) for a family of indeterminates. The symmetric formal power series \(s_{k}\) are defined as before by

\[
s _ {k} = \sum_ {\mathrm{H} \in \mathfrak {P} _ {k}} \prod_ {i \in \mathrm{H}} X _ {i} \quad (k \geqslant 1),\tag{16}
\]

where \(P_{k}\) is the set of all k-element subsets of I. We also write

\[
p _ {k} = \sum_ {i \in I} X _ {i} ^ {k} \quad (k \geqslant 1).\tag{17}
\]

This is a symmetric formal power series which is homogeneous of degree k.

Lemma 4 (« Newton's relations »). --- For every integer \(d \geqslant 1\) we have

\[
p _ {d} = \sum_ {k = 1} ^ {d - 1} (- 1) ^ {k - 1} s _ {k} p _ {d - k} + (- 1) ^ {d + 1} d s _ {d}.\tag{18}
\]

Let us define a continuous derivation \(\Delta\) in \(\mathbf{A}[[\mathbf{X}]]\) by \(\Delta(u) = \sum_{n \geq 0} nu_n\) , where for every \(u\) in \(\mathbf{A}[[\mathbf{X}]]\) , \(u_n\) is the homogeneous component of degree \(n\) of \(u\) . By (16) and Prop. 2 of IV, p.~27 we have

\[
1 + \sum_ {k \geqslant 1} s _ {k} = \prod_ {i \in I} (1 + X _ {i}).\tag{19}
\]

By Prop. 9 of IV, p.~34 we thus have

\[
\left(\sum_ {k \geqslant 1} k s _ {k}\right) \cdot \left(1 + \sum_ {k \geqslant 1} s _ {k}\right) ^ {- 1} = \sum_ {i \in I} \Delta \left(\mathrm{X} _ {i}\right) / \left(1 + \mathrm{X} _ {i}\right).\tag{20}
\]

We have \(\Delta(X_i) = X_i\) and \(X_i / (1 + X_i) = \sum_{k \geq 1} (-1)^{k-1} X_i^k\) . The right-hand side of (20) is thus equal to \(\sum_{k \geq 1} (-1)^{k-1} p_k\) . Hence by (20) we have

\[
\sum_ {k \geqslant 1} k s _ {k} = \left(1 + \sum_ {k \geqslant 1} s _ {k}\right) \cdot \left(\sum_ {k \geqslant 1} (- 1) ^ {k - 1} p _ {k}\right)
\]

and now Lemma 4 follows by a comparison of homogeneous components of degree \(d\) .

Remark. --- With the notation of the above proof we have

\[
\Delta u = \sum_ {i \in I} X _ {i} \cdot D _ {i} (u)
\]

(IV, p.~33, Cor. 1). In other words, Euler's relation (IV, p.~8, Prop 6) extends to formal power series: if \(u \in \mathbf{A}[\mathbf{X}]\) is homogeneous of degree \(n\) , then

\[
n \cdot u = \sum_ {i \in I} X _ {i} \cdot D _ {i} (u).\tag{21}
\]

When I is finite, consisting of \(n\) elements, we have \(s_k = 0\) for \(k > n\) . Then Newton's relations may be written

\[
\begin{array}{l} p _ {2} = s _ {1} p _ {1} - 2 s _ {2} \\ p _ {3} = s _ {1} p _ {2} - s _ {2} p _ {1} + 3 s _ {3} \\ \dots \dots \\ p _ {n - 1} = s _ {1} p _ {n - 2} - s _ {2} p _ {n - 3} + \dots + (- 1) ^ {n - 1} s _ {n - 2} p _ {1} + (- 1) ^ {n} (n - 1) s _ {n - 1} \\ p _ {n} = s _ {1} p _ {n - 1} - s _ {2} p _ {n - 2} + \dots + (- 1) ^ {n} s _ {n - 1} p _ {1} + (- 1) ^ {n + 1} n s _ {n} \end{array}
\]

and

\[
p _ {k} = s _ {1} p _ {k - 1} - s _ {2} p _ {k - 2} + \dots + (- 1) ^ {n + 1} s _ {n} p _ {k - n} \quad (\text { for } k > n).\tag{22}
\]

The first n of the above relations hold for any I; we find for example

\[
p _ {1} = s _ {1}, \quad p _ {2} = s _ {1} ^ {2} - 2 s _ {2}, \quad p _ {3} = s _ {1} ^ {3} - 3 s _ {1} s _ {2} + 3 s _ {3}.
\]

More generally, let \(\mathbf{S} = (\mathbf{S}_n)_{n\geqslant 1}\) be a family of indeterminates. Let us define the polynomials \(\mathbf{P}_d\in \mathbf{Z}[\mathbf{S}_1,\dots ,\mathbf{S}_d]\) recursively by \(\mathbf{P}_1 = \mathbf{S}_1\) and

\[
\mathrm{P} _ {d} = \sum_ {k = 1} ^ {d - 1} (- 1) ^ {k - 1} \mathrm{S} _ {k} \mathrm{P} _ {d - k} + (- 1) ^ {d + 1} d \mathrm{S} _ {d} \quad (d \geqslant 2).
\]

Then we have the « universal formulae » \(p_d = \mathrm{P}_d(s_1, \ldots, s_d)\) holding over any ring A and family X of indeterminates.

Let \(\mathbf{P} = (\mathrm{P}_k)_{k\geqslant 1}\) be a sequence of indeterminates. Since \(p_k\) is homogeneous of degree \(k\) in \(\mathbf{A}[[\mathbf{X}]]\) , there exists precisely one continuous A-algebra homomorphism

\[
\lambda_ {I}: \mathbf {A} [ [ \mathbf {P} ] ] \rightarrow \mathbf {A} [ [ \mathbf {X} ] ] ^ {\text { sym }}
\]

such that \(\lambda_{I}(\mathbf{P}_{k}) = p_{k}\) for every integer \(k \geqslant 1\) (IV, p.~28). If we assign to \(\mathbf{P}_{k}\) the weight \(k\) , then \(\lambda_{I}\) transforms an isobaric polynomial of weight \(n\) in the \(\mathbf{P}_{k}\) into a formal power series which is homogeneous of degree \(n\) in the \(\mathbf{X}_{i}\) .

PROPOSITION 3. --- a) If I is a finite set with n elements and \(n! \cdot 1\) is invertible in A, then \(\lambda_{I}\) induces a bicontinuous isomorphism of \(A[[P_{1}, ..., P_{n}]]\) onto \(A[[X]]^{sym}\) .

\begin{enumerate}
\def\labelenumi{\alph{enumi})}
\setcounter{enumi}{1}
\tightlist
\item
  If I is infinite and A is a Q-algebra, then \(\lambda_{I}\) is a bicontinuous isomorphism of A{[}{[}P{]}{]} onto A{[}{[}X{]}{]} \(^{\text{sym}}\) .
\end{enumerate}

We shall treat the case a) only, the case b) being quite similar.

By Th. 2 (IV, p.~68) we can identify A{[}{[}X{]}{]} \(^{sym}\) with the algebra of formal power series A{[}{[}S \(_1\) , \ldots, S \(_n\) {]}{]}, with S \(_k\) corresponding to s \(_k\) . By Lemma 4 of IV, p.~70, there exist formal power series g \(_1\) , \ldots, g \(_n\) of order \(\geqslant 2\) in the indeterminates s \(_1\) , \ldots, s \(_n\) such that

\[
p _ {k} = (- 1) ^ {k + 1} k s _ {k} + g _ {k} (s _ {1}, \dots , s _ {n}) \quad (1 \leqslant k \leqslant n).
\]

Since \(k!\) . 1 is invertible in A, Lemma 2 of IV, p.~35 proves the existence of an automorphism T of the topological A-algebra \(\mathbf{A}[[\mathbf{X}]]^{\mathrm{sym}}\) which maps \(s_k\) to \(p_k\) for \(1 \leqslant k \leqslant n\) . Now Prop. 3, a) is an immediate consequence.

COROLLARY. --- Let \(\xi_1, \ldots, \xi_n, \eta_1, \ldots, \eta_n\) be elements of A and suppose that A is an integral domain.

\begin{enumerate}
\def\labelenumi{\alph{enumi})}
\item
  If \(s_k(\xi_1, \ldots, \xi_n) = s_k(\eta_1, \ldots, \eta_n)\) for \(1 \leqslant k \leqslant n\) , then there exists a permutation \(\sigma \in \mathfrak{S}_n\) such that \(\eta_i = \xi_{\sigma(i)}\) for \(1 \leqslant i \leqslant n\) .
\item
  Suppose that \(n!\) . \(1 \neq 0\) in A and
\end{enumerate}

\[
\xi_ {1} ^ {k} + \dots + \xi_ {n} ^ {k} = \eta_ {1} ^ {k} + \dots + \eta_ {n} ^ {k}\tag{23}
\]

for \(1 \leqslant k \leqslant n\) . Then there exists a permutation \(\sigma \in \mathfrak{S}_n\) such that \(\eta_i = \xi_{\sigma(i)}\) for \(1 \leqslant i \leqslant n\) .

Under the hypotheses \(a)\) we have \(\prod_{i=1}^{n} (\mathbf{X} - \xi_i) = \prod_{i=1}^{n} (\mathbf{X} - \eta_i)\) . If we substitute \(\eta_n\) for \(\mathbf{X}\) , we find that \(\prod_{i=1}^{n} (\eta_n - \xi_i) = 0\) and since \(\mathbf{A}\) is an integral domain, there is an integer \(\sigma(n)\) such that \(1 \leqslant \sigma(n) \leqslant n\) and \(\eta_n = \xi_{\sigma(n)}\) . Now the assertion \(a)\) follows easily by induction because \(\mathbf{A}[\mathbf{X}]\) is an integral domain.

Under the hypotheses of \(b)\) there exist by Prop. 3 polynomials \(\Pi_1, \ldots, \Pi_n\) in \(n\) indeterminates, with coefficients in the field of fractions of \(A\) such that \(s_k = \Pi_k(p_1, \ldots, p_n)\) for \(1 \leqslant k \leqslant n\) . Now the relation (23) implies

\[
s _ {k} (\xi_ {1}, \dots , \xi_ {n}) = s _ {k} (\eta_ {1}, \dots , \eta_ {n})
\]

for \(1 \leqslant k \leqslant n\) and so \(b)\) follows from \(a)\) .

\subsection{Symmetric functions in the roots of a polynomial}

Consider a monic polynomial of degree \(n\) , with coefficients in \(A\) :

\[
f = \mathbf {X} ^ {n} + a _ {1} \mathbf {X} ^ {n - 1} + \dots + a _ {n - 1} \mathbf {X} + a _ {n}.
\]

We define an associative, commutative and unital A-algebra \(\mathbf{E}_f\) with the generators \(x_{1},\ldots ,x_{n}\) and relations

\[
\sum_ {i _ {1} <   \dots <   i _ {k}} x _ {i _ {1}} \dots x _ {i _ {k}} = (- 1) ^ {k} a _ {k} \quad (1 \leqslant k \leqslant n).\tag{24}
\]

More precisely, we have

\[
\mathrm{E} _ {f} = \mathrm{A} [ \mathrm{X} _ {1}, \dots , \mathrm{X} _ {n} ] / \mathfrak {a}
\]

where the ideal \(\mathfrak{a}\) is generated by the polynomials \(s_k + (-1)^{k+1} a_k\) for \(1 \leqslant k \leqslant n\) , and \(x_i\) is the residue class of \(\mathbf{X}_i\) mod \(\mathfrak{a}\) for \(1 \leqslant i \leqslant n\) . The relation (24) is also equivalent to \(f(\mathbf{X}) = \prod_{i=1}^{n} (\mathbf{X} - x_i)\) . When there is any risk of ambiguity, we write \(x_{1,f}, \ldots, x_{n,f}\) in place of \(x_1, \ldots, x_n\) .

PROPOSITION 4. --- Let B be a commutative ring, \(\rho\) a homomorphism of A into B and \(\xi_1, \ldots, \xi_n\) elements of B. Suppose that the relation \( {}^{\rho}f(X) = \prod_{i=1}^{n}(X - \xi_i)\) holds in B{[}X{]}. Then there exists one and only one ring homomorphism \(u: E_f \to B\) such that \(\rho(a) = u(a, 1)\) for all \(a \in A\) and \(u(x_i) = \xi_i\) for \(1 \leqslant i \leqslant n\) .

We consider B as associative, commutative and unital A-algebra by means of \(\rho\) . Then the relation \({}^{\rho}f(X) = \prod_{i=1}^{n}(X - \xi_i)\) may also be written in the form

\[
\sum_ {i _ {1} <   \dots <   i _ {k}} \xi_ {i _ {1}} \dots \xi_ {i _ {k}} = (- 1) ^ {k} a _ {k} \cdot 1 \quad (1 \leqslant k \leqslant n)
\]

in B. Since the relations (24) define a presentation of \(\mathbf{E}_f\) , Prop. 4 follows.

Prop. 4 justifies the name « universal decomposition algebra of f» for \(E_{f}\) . The relation \(f(\mathbf{X}) = \prod_{i=1}^{n} (\mathbf{X} - x_{i,f})\) is called the « universal decomposition of f». Let \(\sigma \in S_{n}\) be a permutation; since \(f(\mathbf{X}) = \prod_{i=1}^{n} (\mathbf{X} - x_{\sigma(i),f})\) , there exists an automorphism \(t_{\sigma}\) of the A-algebra \(E_{f}\) characterized by \(t_{\sigma}(x_{i,f}) = x_{\sigma(i),f}\) for \(1 \leqslant i \leqslant n\) . We have \(t_{\sigma\tau} = t_{\sigma} \circ t_{\tau}\) for \(\sigma, \tau\) in \(S_{n}\) , and hence obtain an action of the group \(S_{n}\) on the A-algebra \(E_{f}\) .

PROPOSITION 5. --- In the universal decomposition algebra \(E_{f}\) the family of monomials \(x_{1}^{\nu(1)}\ldots x_{n}^{\nu(n)}\) such that \(0 \leqslant \nu(i) < i\) for \(1 \leqslant i \leqslant n\) is a basis of the A-module \(E_{f}\) . In particular \(E_{f}\) is a free A-module of rank \(n!\) .

Write \(\mathbf{B} = \mathbf{A}[X_1, \ldots, X_n]\) and \(\mathbf{C} = \mathbf{A}[X_1, \ldots, X_n]^{\mathrm{sym}}\) . By Th. 1 (IV, p.~62) we have \(\mathbf{C} = \mathbf{A}[s_1, \ldots, s_n]\) and \(s_1, \ldots, s_n\) are algebraically independent over \(\mathbf{A}\) . The polynomials without constant term in \(s_1, \ldots, s_n\) form an ideal \(\mathbf{C}^+\) in \(\mathbf{C}\) , supplementary to \(\mathbf{A}\) and generated by \(s_1, \ldots, s_n\) . Let \(\mathfrak{c}\) be the ideal of \(\mathbf{C}\) generated by \(s_1 + a_1\) , \(s_2 - a_2\) , \ldots, \(s_n + (-1)^{n+1}a_n\) . There exists an \(\mathbf{A}\) -algebra automorphism of \(\mathbf{C}\) which maps \(s_k\) to \(s_k + (-1)^{k+1}a_k\) for \(1 \leqslant k \leqslant n\) , and hence maps \(\mathbf{C}^+\) onto \(\mathfrak{c}\) ; therefore we have \(\mathbf{C} = \mathbf{A} \oplus \mathfrak{c}\) . Moreover Th. 1, c) of IV, p.~62 shows that

\[
\mathbf {B} = \oplus \mathbf {C X} ^ {\nu}
\]

where S is the set of all \(\nu \in \mathbf{N}^n\) such that \(0 \leqslant \nu(i) < i\) for \(1 \leqslant i \leqslant n\) . The ideal \(\mathfrak{a}\) of B is generated by \(\mathfrak{c}\) , whence \(\mathfrak{a} = \mathbf{B}\mathfrak{c} = \bigoplus_{\nu \in \mathbb{S}} \mathfrak{c} \cdot \mathbf{X}^\nu\) . Since \(\mathbb{C} = \mathbb{A} \oplus \mathfrak{c}\) , we find

\[
\mathbf {B} = \mathfrak {a} \oplus \bigoplus_ {\nu \in S} \mathbf {A X} ^ {\nu},
\]

whence Prop. 5, because \(\mathbf{E}_f = \mathbf{B} / \mathfrak{a}\) .

COROLLARY. --- The canonical homomorphism of A into the universal decomposition algebra of the monic polynomial \(f \in A[X]\) is injective.

For the unit element of \(E_{f}\) forms part of a basis of the A-module \(E_{f}\) .

PROPOSITION 6. --- Let \(f \in \mathbf{A}[\mathbf{X}]\) be a monic polynomial of degree \(n\) , and let \(P\) be a symmetric polynomial in \(\mathbf{X}_1, \ldots, \mathbf{X}_n\) with coefficients in \(A\) . Then there exists precisely one element \(a\) of \(A\) with the following property:

(FS) For any ring homomorphism \(\rho: \mathbf{A} \to \mathbf{B}\) and decomposition \({}^{\rho} f(\mathbf{X}) = \prod_{i=1}^{n} (\mathbf{X} - \xi_i)\) in \(\mathbf{B}[\mathbf{X}]\) we have \(\rho(a) = \mathrm{P}(\xi_1, \ldots, \xi_n)\) .

Write \(f = \mathbf{X}^n + \sum_{k=1}^{n} a_k \mathbf{X}^{n-k}\) , then by Th. 1, IV, p.~62 there exists a polynomial

\(\Pi\) in \(n\) indeterminates with coefficients in \(\mathbf{A}\) , such that \(\mathbf{P} = \Pi(s_1, \ldots, s_n)\) . Put \(a = \Pi(-a_1, a_2, \ldots, (-1)^n a_n)\) . Under the hypothesis (FS) we have

\[
s _ {k} (\xi_ {1}, \dots , \xi_ {n}) = (- 1) ^ {k} \rho (a _ {k})
\]

whence

\[
\begin{array}{r l} \rho (a) & = \Pi (- \rho (a _ {1}), \rho (a _ {2}), \dots , (- 1) ^ {n} \rho (a _ {n})) \\ & = \Pi (s _ {1} (\xi_ {1}, \dots , \xi_ {n}), \dots , s _ {n} (\xi_ {1}, \dots , \xi_ {n})) \\ & = P (\xi_ {1}, \dots , \xi_ {n}). \end{array}
\]

This proves the existence of an element a satisfying (FS). The uniqueness of a follows from the Cor. to Prop. 5, because we have \(a \cdot 1 = P(x_{1,f}, \ldots, x_{n,f})\) in the universal decomposition algebra \(\mathbf{E}_f\) .

With the notation of Prop. 6 one sometimes writes \(a = \mathrm{P}^{\#}(f)\) . Here are some examples.

Examples. --- 1) If \(\mathrm{P} = s_k\) , then \(\mathrm{P}^{\#}(f) = (-1)^k a_k\) .

2) Let \(g\) be a polynomial in \(\mathbf{A}[\mathbf{X}]\) and put

\[
\mathrm{P} \left(\mathrm{X} _ {1}, \dots , \mathrm{X} _ {n}\right) = g \left(\mathrm{X} _ {1}\right) \dots g \left(\mathrm{X} _ {n}\right).
\]

Then \(\mathrm{P}^{\#}(f)\) is just the resultant \(\operatorname{res}(f, g)\) , by Cor. 1 of IV, p.~80.

\begin{enumerate}
\def\labelenumi{\arabic{enumi})}
\setcounter{enumi}{2}
\item
  Put \(\Delta(X_1, \ldots, X_n) = \prod_{i < j} (X_i - X_j)^2\) , then \(\Delta^{\#}(f)\) is just the discriminant of the monic polynomial \(f\) (IV, p.~82, formula (46)).
\item
  Put \(\mathrm{P}(\mathrm{X}_{1},\ldots,\mathrm{X}_{n})=\mathrm{X}_{1}^{k}+\cdots+\mathrm{X}_{n}^{k}\) ; moreover, define the algebra \(\mathrm{E}=\mathrm{A}[\mathrm{X}]/(f)\) and write x for the image of X in E. We recall that the A-module E is free, with basis \((1,x,\ldots,x^{n-1})\) (IV, p.~11, Cor.). Let us show
\end{enumerate}

\[
\operatorname{Tr} _ {\mathrm{E/A}} (x ^ {k}) = \mathrm{P} ^ {\#} (f).\tag{25}
\]

Put \(\pi_k = \mathrm{Tr}_{\mathrm{E / A}}(x^k)\) for every integer \(k \geqslant 1\) . Bearing in mind Newton's relations (IV, p.~70) we need only establish the relations

\begin{enumerate}
\def\labelenumi{(\arabic{enumi})}
\setcounter{enumi}{25}
\tightlist
\item
\end{enumerate}

\[
\begin{array}{l l} \pi_ {k} + a _ {1} \pi_ {k - 1} + \dots + a _ {k - 1} \pi_ {1} + k a _ {k} = 0 & \text { for } \quad 1 \leqslant k \leqslant n \\ \pi_ {k} + a _ {1} \pi_ {k - 1} + \dots + a _ {n - 1} \pi_ {k - n + 1} + a _ {n} \pi_ {k - n} = 0 & \text { for } \quad k > n \end{array}\tag{27}
\]

(which we shall also call « Newton's relations »). The relation (27) is clear, because the left-hand side is the trace of

\[
x ^ {k - n} \left(x ^ {n} + a _ {1} x ^ {n - 1} + \dots + a _ {n - 1} x + a _ {n}\right) = 0.
\]

Suppose that \(1 \leqslant k \leqslant n\) and put

\[
y = x ^ {k} + a _ {1} x ^ {k - 1} + \dots + a _ {k - 1} x + a _ {k} \cdot 1;
\]

let \(\mathbf{M} = (m_{ij})\) be the matrix of the linear mapping \(u\mapsto yu\) in E relative to the basis \((x^{i})_{0\leqslant i\leqslant n-1}\) . We easily obtain the relations

\[
\begin{array}{l l} m _ {i i} = a _ {k} & \text { for } \quad 0 \leqslant i <   n - k \\ m _ {i i} = 0 & \text { for } \quad n - k \leqslant i <   n, \end{array}
\]

whence

\[
\operatorname{Tr} _ {\mathrm{E/A}} (y) = \sum_ {i = 0} ^ {n - 1} m _ {i i} = (n - k) a _ {k}.
\]

Moreover we have

\[
\operatorname{Tr} _ {\mathrm{E/A}} (y) = \pi_ {k} + a _ {1} \pi_ {k - 1} + \dots + a _ {k - 1} \pi_ {1} + n a _ {k},
\]

whence formula (26) follows.

\subsection{The resultant}

In this No.~we assume as given two positive integers \(p, q\) and two polynomials \(f, g\) in \(\mathbf{A}[\mathbf{X}]\) of the form

\[
\begin{array}{r l} & {f = t _ {p} \mathbf {X} ^ {p} + t _ {p - 1} \mathbf {X} ^ {p - 1} + \dots + t _ {0}} \\ & {g = u _ {q} \mathbf {X} ^ {q} + u _ {q - 1} \mathbf {X} ^ {q - 1} + \dots + u _ {0}} \end{array}
\]

such that \(\deg f \leqslant p\) , \(\deg g \leqslant q\) . For every integer \(n \geqslant 0\) we denote by \(S_n\) the sub-A-module of A{[}X{]} consisting of all polynomials of degree \(< n\) ; it has the family \((X^i)_{0 \leqslant i < n}\) as basis, and hence is of rank \(n\) .

We provide \(\mathbf{S}_q\times \mathbf{S}_p\) with the basis

\[
\mathrm{B} _ {1} = \left(\left(\mathrm{X} ^ {q - 1}, 0\right), \dots , (\mathrm{X}, 0), (1, 0), \left(0, \mathrm{X} ^ {p - 1}\right), \dots , (0, \mathrm{X}), (0, 1)\right)
\]

and \(\mathbf{S}_{p + q}\) with the basis

\[
\mathbf {B} _ {2} = \left(\mathbf {X} ^ {p + q - 1}, \dots , \mathbf {X}, 1\right).
\]

We define a linear mapping \(\varphi : S_q \times S_p \to S_{p + q}\) by

\[
\varphi (u, v) = u f + v g
\]

and we denote by \(\mathbf{M}(f,g,p,q)\) the matrix of \(\varphi\) with respect to the bases \(\mathbf{B}_1\) and \(\mathbf{B}_2\) . This is a square matrix of order \(p + q\) , indexed by the set \(\{0,1,\dots,p + q - 1\}\) . Its elements \(a_{ij}\) are given by the rules:

\[
\begin{array}{l l} a) & a _ {i j} = t _ {p - i + j} \quad \text {for} \quad 0 \leqslant j \leqslant q - 1, \\ b) & a _ {i j} = u _ {j - i} \quad \text {for} \quad q \leqslant j \leqslant p + q - 1, \end{array}
\]

where \(t_k\) is taken to be 0 if \(k \notin [0, p]\) and \(u_k = 0\) if \(k \notin [0, q]\) .

For example, when \(p = 2\) and \(q = 3\) , we have the matrix

\[
\left( \begin{array}{c c c c c} t _ {2} & 0 & 0 & u _ {3} & 0 \\ t _ {1} & t _ {2} & 0 & u _ {2} & u _ {3} \\ t _ {0} & t _ {1} & t _ {2} & u _ {1} & u _ {2} \\ 0 & t _ {0} & t _ {1} & u _ {0} & u _ {1} \\ 0 & 0 & t _ {0} & 0 & u _ {0} \end{array} \right).
\]

DEFINITION 1. --- With the above notation the determinant of the matrix \(\mathbf{M}(f,g,p,q)\) is called the resultant of the pair \((f,g)\) for the degrees p and q, or simply the resultant of f and g if \(p = \deg f\) and \(q = \deg g\) .

The resultant is denoted by \(\operatorname{res}_{p,q}(f,g)\) or simply \(\operatorname{res}(f,g)\) when \(p=\deg f\) , \(q=\deg g\) .

Examples. --- 1) Given \(\lambda\) , \(\mu\) in A, we have the formulae

\[
\begin{array}{l} \operatorname{res} _ {p, 0} (f, \lambda) = \lambda^ {p}, \quad \operatorname{res} _ {0, q} (\mu , g) = \mu^ {q} \\ \operatorname{res} _ {p, 1} (f, \lambda) = \lambda^ {p} t _ {p}, \quad \operatorname{res} _ {1, q} (\mu , g) = (- 1) ^ {q} \mu^ {q} u _ {q}, \end{array}
\]

whose proof is immediate.

\begin{enumerate}
\def\labelenumi{\arabic{enumi})}
\setcounter{enumi}{1}
\tightlist
\item
  When \(p = q = 1\) , we have
\end{enumerate}

\[
\operatorname{res} _ {1, 1} \left(t _ {1} \mathbf {X} + t _ {0}, u _ {1} \mathbf {X} + u _ {0}\right) = t _ {1} u _ {0} - t _ {0} u _ {1}.
\]

Remarks. --- 1) The matrix \(\mathbf{M}(g, f, q, p)\) is obtained from \(\mathbf{M}(f, g, p, q)\) by \(pq\) transpositions of columns, whence

\[
\operatorname{res} _ {q, p} (g, f) = (- 1) ^ {p q} \operatorname{res} _ {p, q} (f, g).
\]

\begin{enumerate}
\def\labelenumi{\arabic{enumi})}
\setcounter{enumi}{1}
\tightlist
\item
  Let \(\rho: \mathbf{A} \to \mathbf{B}\) be a ring homomorphism. Def. 1 implies directly the formula
\end{enumerate}

\[
\operatorname{res} _ {p, q} \left({}^ {\rho} f, {}^ {\rho} g\right) = \rho \left(\operatorname{res} _ {p, q} (f, g)\right).
\]

\begin{enumerate}
\def\labelenumi{\arabic{enumi})}
\setcounter{enumi}{2}
\tightlist
\item
  Given \(\lambda, \mu\) in A, we have
\end{enumerate}

\[
\operatorname{res} _ {p, q} (\lambda f, \mu g) = \lambda^ {q} \mu^ {p} \operatorname{res} _ {p, q} (f, g).\tag{28}
\]

\begin{enumerate}
\def\labelenumi{\arabic{enumi})}
\setcounter{enumi}{3}
\tightlist
\item
  Suppose that \(p + q \geqslant 1\) . By III, p.~532, formula (28), the image of \(\varphi\) contains the constant polynomial \(\operatorname{res}_{p,q}(f, g)\) . Hence there exist a pair of polynomials \((u, v)\) , with \(u \in S_q\) , \(v \in S_p\) such that
\end{enumerate}

\[
\operatorname{res} _ {p, q} (f, g) = u f + v g,
\]

whence

\[
\operatorname{res} _ {p, q} (f, g) \in \mathrm{A} \cap (f, g).
\]

This pair \((u, v)\) is unique when \(\operatorname{res}_{p,q}(f, g)\) is cancellable in \(\mathbf{A}\) : for then \(\varphi\) is injective (III, p.~524, Prop. 3).

\begin{enumerate}
\def\labelenumi{\arabic{enumi})}
\setcounter{enumi}{4}
\tightlist
\item
  Suppose that \(p \geqslant q\) , and let \(h \in \mathbf{A}[\mathbf{X}]\) be a polynomial of degree \(\leqslant p - q\) . Let us show that
\end{enumerate}

\[
\operatorname{res} _ {p, q} (f, g) = \operatorname{res} _ {p, q} (f + g h, g).\tag{29}
\]

For if we write \(\omega(u, v) = (u, uh + v)\) for \((u, v) \in S_q \times S_p\) , then \(\omega\) is an automorphism of the A-module \(S_q \times S_p\) and we have

\[
\omega^ {- 1} (u, v) = (u, - u h + v).
\]

The matrix of \(\omega\) with respect to the basis \(B_{1}\) is lower triangular and its diagonal elements are equal to 1. On the other hand, \(\varphi \circ \omega\) maps \((u, v)\) to \(u(f + gh) + vg\) . Now formula (29) means that the matrices representing \(\varphi\) and \(\varphi \circ \omega\) have the same determinant and this follows from the relation det \(\omega = 1\) .

Assume that \(f\) is monic of degree \(p\) ; put \(E = A[X]/(f)\) and denote by \(x\) the canonical image of \(X\) in \(E\) . We know (IV, p.~11) that \(E\) is a free \(A\) -module with basis \((1, x, ..., x^{p-1})\) . We can thus define the norm \(N_{E/A}(u)\) of any element \(u\) of \(E\) (III, p.~543, Def. 2).

PROPOSITION 7. --- Assume that \(f\) is monic of degree \(p\) ; with the above notations we have \(^{1}\)

\[
\operatorname{res} _ {p, q} (f, g) = \mathrm{N} _ {\mathrm{E/A}} (g (x)).\tag{30}
\]

Let us define an A-linear mapping \(\theta\) of \(S_q \times S_p\) into \(S_{p+q}\) by \(\theta(u, v) = uf + v\) . Then \(\theta\) maps the basis \(B_1\) of \(S_q \times S_p\) into the sequence

\[
(f \mathbf {X} ^ {q - 1}, \dots , f \mathbf {X}, f, \mathbf {X} ^ {p - 1}, \dots , \mathbf {X}, 1)
\]

of elements of \(S_{p+q}\) ; hence the matrix \(M_{\theta}\) of \(\theta\) with respect to the bases \(B_{1}\) , \(B_{2}\) is lower triangular and its diagonal elements are equal to 1, whence det \(M_{\theta}=1\) .

It follows that \(\theta\) is bijective and \(\operatorname{res}_{p,q}(f,g)\) is equal to the determinant of the endomorphism \(\varphi' = \varphi \circ \theta^{-1}\) of \(S_{p + q}\) . Explicitly, we have

\[
\varphi^ {\prime} (u f + v) = u f + v g\tag{31}
\]

for any pair \((u,v)\) in \(S_{q}\times S_{p}\) . Now we have \(A[X]=S_{p+q}+(f)\) and \(fS_{q}=S_{p+q}\cap(f)\) , hence the canonical injection of \(S_{p+q}\) into A{[}X{]} defines by the passage to quotients an isomorphism \(\gamma\) of \(S_{p+q}/fS_{q}\) onto E. Let \(\psi\) be the multiplication by \(g(x)\) in E. Formula (31) shows that \(\varphi'\) induces the identity on \(fS_{q}\) and \(\gamma^{-1}\psi\gamma\) on \(S_{p+q}/fS_{q}\) . We thus have det \(\varphi'=\det\psi\) , whence (30) follows, because det \(\varphi'=\operatorname{res}_{p,q}(f,g)\) and det \(\psi=N_{E/A}(g(x))\) , by definition.

COROLLARY 1. --- Let \(f \in \mathbf{A}[\mathbf{X}]\) be a monic polynomial; for every polynomial \(g \in \mathbf{A}[\mathbf{X}]\) the following conditions are equivalent:

\begin{enumerate}
\def\labelenumi{(\roman{enumi})}
\item
  res \((f,g)\) is invertible in A;
\item
  there exist polynomials \(u, v\) in \(\mathbf{A}[\mathbf{X}]\) such that \(uf + vg = 1\) ;
\item
  \(g(x)\) is invertible in the algebra \(\mathbf{A}[\mathbf{X}] / (f)\) .
\end{enumerate}

The equivalence of (i) and (iii) follows from Prop. 3 of III, p.~545 and Prop. 7; that of (ii) and (iii) follows trivially.

COROLLARY 2. --- Suppose that A is a field and let \(f, g \in A[X]\) , then the following conditions are equivalent when \(f, g\) are non-zero:

\begin{enumerate}
\def\labelenumi{(\roman{enumi})}
\item
  \(\operatorname{res}(f, g) \neq 0\) ;
\item
  the polynomials \(f\) and \(g\) are relatively prime in \(\mathbf{A}[\mathbf{X}]\) ;
\end{enumerate}

(iii) for any extension L of A the polynomials f and g have no common root in L.

We can reduce at once to the case where \(f\) is monic (IV, p.~77, remark 3). The equivalence of (i) and (ii) is just a translation of Cor. 1, by Def. 1 of IV, p.~12; the equivalence of (ii) and (iii) is just Cor. 7 of IV, p.~13.

COROLLARY 3. --- For every \(\lambda \in A\) we have

\[
\operatorname{res} _ {p, 1} (f, \lambda - X) = f (\lambda), \quad \operatorname{res} _ {1, q} (X - \lambda , g) = g (\lambda).\tag{32}
\]

When \(f(\mathbf{X}) = \mathbf{X} - \lambda\) , the algebra E is equal to A and we have \(x = \lambda\) ; the second formula (32) now follows from Prop. 7 (IV, p.~77). By remarks 1 and 3 (IV, p.~76, 77) we conclude

\[
\operatorname{res} _ {p, 1} (f, \lambda - \mathbf {X}) = (- 1) ^ {p} \operatorname{res} _ {1, p} (\lambda - \mathbf {X}, f) = (- 1) ^ {p + p} \operatorname{res} _ {1, p} (\mathbf {X} - \lambda , f) = f (\lambda).
\]

Suppose now that f and g are monic. We denote by F the A-algebra \(\mathrm{A}[\mathrm{X},\mathrm{Y}]/(f(\mathrm{X}),g(\mathrm{Y}))\) and denote by x (resp. y) the canonical image of X (resp. Y) in F.

PROPOSITION 8. --- Suppose that \(f\) and \(g\) are monic of degrees \(p\) and \(q\) respectively. With the above notation, the A-module \(F\) is free with basis \((x^i y^j)_{0 \leq i < p, 0 \leq j < q}\) and we have

\[
\operatorname{res} (f, g) = \mathrm{N} _ {\mathrm{F/A}} (x - y).\tag{33}
\]

Put \(E = A[X]/(f)\) and \(E' = A[Y]/(g)\) . By II, p.~253, Cor. 1, the homomorphism \(\sigma\) of \(E \otimes E'\) into F derived from the canonical homomorphism \(A[X] \otimes A[Y] \to A[X,Y]\) is bijective; this proves the assertion about the basis of F. We now identify E with its image in F under \(\sigma\) . Then the E-algebra homomorphism of \(E[Y]/(g(Y))\) into F which maps Y to y is an isomorphism.

By the transitivity of the norm (III, p.~546) we have

\[
\mathrm{N} _ {\mathrm{F/A}} (x - y) = \mathrm{N} _ {\mathrm{E/A}} (\mathrm{N} _ {\mathrm{F/E}} (x - y)).\tag{34}
\]

By Prop. 7 (IV, p.~77), \(\mathbf{N}_{\mathrm{F / E}}(x - y)\) is the resultant of the polynomials \(g(\mathbf{Y})\) and \(x - \mathbf{Y}\) in E{[}Y{]}, hence equal to \(g(x)\) (IV, p.~78, Cor. 3). By formula (34) and Prop. 7 (IV, p.~77) we thus have

\[
\mathrm{N} _ {\mathrm{F/A}} (x - y) = \mathrm{N} _ {\mathrm{E/A}} (g (x)) = \operatorname{res} (f, g).
\]

PROPOSITION 9. --- Let \(p_1\) and \(q_1\) be positive integers and \(f_1, g_1\) polynomials in A{[}X{]} such that \(\deg f_1 \leqslant p_1\) , \(\deg g_1 \leqslant q_1\) ; then we have

\begin{enumerate}
\def\labelenumi{(\arabic{enumi})}
\setcounter{enumi}{34}
\tightlist
\item
\end{enumerate}

\[
\operatorname{res} _ {p, q + q _ {1}} (f, g g _ {1}) = \operatorname{res} _ {p, q} (f, g) \cdot \operatorname{res} _ {p, q _ {1}} (f, g _ {1})\tag{36}
\]

\[
\operatorname{res} _ {p + p _ {1}, q} (f f _ {1}, g) = \operatorname{res} _ {p, q} (f, g) \cdot \operatorname{res} _ {p _ {1}, q} (f _ {1}, g).
\]

We have \(\operatorname{res}_{p,q}(f,g)=(-1)^{pq}\operatorname{res}_{q,p}(g,f)\) (IV, p.~76); so it is enough to prove (35). Likewise, Remark 3 (loc. cit.) shows that if formula (35) is established for a polynomial f, then it holds for all polynomials of the form \(\lambda f\) , where \(\lambda\in A\) . Finally, when f is monic of degree p, formula (35) follows from Prop. 7 (IV, p.~77) by the formula \(\mathrm{N}_{\mathrm{E}/\mathrm{A}}(ab)=\mathrm{N}_{\mathrm{E}/\mathrm{A}}(a)\cdot\mathrm{N}_{\mathrm{E}/\mathrm{A}}(b)\) . In all we conclude that (35) holds when the coefficient \(t_{p}\) of \(X^{p}\) in f is invertible.

Lemma 5. --- Let t be an element of A. There exists a commutative ring C containing A as subring, a subring B of C containing A, an element \(\tau\) of B invertible in C and a ring homomorphism \(\rho: B \to A\) such that \(\rho(\tau) = t\) and the restriction of \(\rho\) to A is equal to \(Id_{A}\) .

It suffices to take for B the algebra \(\mathbf{A}^{(\mathbf{N})}\) of the monoid N, that is, the polynomial algebra \(A[\tau]\) in one indeterminate \(\tau\) , for C the algebra \(\mathbf{A}^{(\mathbf{Z})}\) of the group Z and for \(\rho\) the homomorphism \(P \mapsto P(t)\) of \(\mathbf{A}[\tau]\) into A.

With the notation of Lemma 5, where we have taken \(t = t_p\) , let us put

\[
\mathbf {F} = \tau \mathbf {X} ^ {p} + t _ {p - 1} \mathbf {X} ^ {p - 1} + \dots + t _ {1} \mathbf {X} + t _ {0}\tag{37}
\]

in B{[}X{]}. The coefficient of \(\mathbf{X}^p\) in F is invertible in C; if we consider F, \(g\) , \(g_1\) as polynomials of C{[}X{]} we thus have

\[
\operatorname{res} _ {p, q + q _ {1}} (\mathrm{F}, g g _ {1}) = \operatorname{res} _ {p, q} (\mathrm{F}, g) \cdot \operatorname{res} _ {p, q _ {1}} (\mathrm{F}, g _ {1})\tag{38}
\]

by what has been said. The resultants are unchanged if we consider F, g and \(g_{1}\) as polynomials in B{[}X{]}. Since \( {}^{\rho}F = f\) , \( {}^{\rho}g = g\) , \( {}^{\rho}g_{1} = g_{1}\) , formula (35) follows from (38) and Remark 2 (IV, p.~77).

COROLLARY 1. --- (i) Let \(\lambda, \alpha_{1}, \ldots, \alpha_{p}\) be elements of A and suppose that \(f(\mathbf{X}) = \lambda(\mathbf{X} - \alpha_{1}) \ldots (\mathbf{X} - \alpha_{p})\) . We have

\[
\operatorname{res} _ {p, q} (f, g) = \lambda^ {q} g (\alpha_ {1}) \dots g (\alpha_ {p}).\tag{39}
\]

\begin{enumerate}
\def\labelenumi{(\roman{enumi})}
\setcounter{enumi}{1}
\tightlist
\item
  Let \(\mu, \beta_1, \ldots, \beta_q\) be elements of \(A\) and suppose further that \(g(X) = \mu(X - \beta_1) \ldots (X - \beta_q)\) . Then we have
\end{enumerate}

\[
\operatorname{res}_{p,q}(f,g) = \lambda^{q}\mu^{p}\prod_{\substack{1\leqslant i\leqslant p\\ 1\leqslant j\leqslant q}}\left(\alpha_{i} - \beta_{j}\right).\tag{40}
\]

Assertion (i) follows directly from the formulae (28), (32) and (36). Now Assertion (ii) follows from (i).

COROLLARY 2. --- For every integer \(r \geqslant 0\) we have

\[
\operatorname{res} _ {p, q + r} (f, g) = t _ {p} ^ {r} \cdot \operatorname{res} _ {p, q} (f, g).\tag{41}
\]

Bearing in mind Example 1 (IV, p.~76) we need only take \(q_{1} = r\) , \(g_{1} = 1\) in (35).

Assume that \(f\) is monic, and let \(\rho: \mathbf{A} \to \mathbf{B}\) be a ring homomorphism and \(\xi_{1}, \ldots, \xi_{p}\) elements of \(\mathbf{B}\) such that we have the decomposition

\[
{ } ^ { p } f ( \mathbf { X } ) = ( \mathbf { X } - \xi _ { 1 } ) \dots ( \mathbf { X } - \xi _ { p } ) .
\]

By Remark 2 (IV, p.~77) and Cor. 1 above we have

\[
\rho (\operatorname{res} (f, g)) = g (\xi_ {1}) \dots g (\xi_ {p}).
\]

This remark applies in particular to the universal decomposition of \(f(\mathrm{IV}, \mathrm{p.~73})\) and since \(\rho\) is then injective, this provides a means of calculating \(\operatorname{res}(f, g)\) .

Example 3. --- Let us prove the formula

\[
\begin{array}{r l} \operatorname{res} _ {2, 2} (a \mathbf {X} ^ {2} + b \mathbf {X} + c, a ^ {\prime} \mathbf {X} ^ {2} + b ^ {\prime} \mathbf {X} + c ^ {\prime}) & = \\ & = (a c ^ {\prime} - c a ^ {\prime}) ^ {2} + (b c ^ {\prime} - c b ^ {\prime}) (b a ^ {\prime} - a b ^ {\prime}). \end{array}\tag{42}
\]

Arguing as for Prop. 9 (IV, p.~79) we see that it is enough to prove the formula when a is invertible. There exists then a decomposition of the form

\[
a \mathbf {X} ^ {2} + b \mathbf {X} + c = a (\mathbf {X} - x) (\mathbf {X} - y)\tag{43}
\]

in B{[}X{]}, where B is an appropriate ring containing A as subring. By Cor. 1 above, the required resultant is equal to

\[
\mathbf {R} = a ^ {2} \left(a ^ {\prime} x ^ {2} + b ^ {\prime} x + c ^ {\prime}\right) \left(a ^ {\prime} y ^ {2} + b ^ {\prime} y + c ^ {\prime}\right).\tag{44}
\]

Now we have

\[
a x + a y = - b, \quad a x y = c
\]

by (43), whence \((ax)^2 + (ay)^2 = b^2 - 2ac\) .

By (44) we have

\[
\begin{array}{r l} \mathrm{R} & = a ^ {\prime 2} (a x y) ^ {2} + a b ^ {\prime 2} (a x y) + a ^ {2} c ^ {\prime 2} + a ^ {\prime} b ^ {\prime} (a x y) (a x + a y) + \\ & \qquad \qquad \qquad \qquad \qquad \qquad \qquad \qquad \qquad \qquad \qquad \qquad + a ^ {\prime} c ^ {\prime} ((a x) ^ {2} + (a y) ^ {2}) + a b ^ {\prime} c ^ {\prime} (a x + a y) \\ & = a ^ {\prime 2} c ^ {2} + a b ^ {\prime 2} c + a ^ {2} c ^ {\prime 2} - a ^ {\prime} b ^ {\prime} c b + a ^ {\prime} c ^ {\prime} (b ^ {2} - 2 a c) - a b ^ {\prime} c ^ {\prime} b \\ & = (a c ^ {\prime} - c a ^ {\prime}) ^ {2} + (a b ^ {\prime} - a ^ {\prime} b) (b ^ {\prime} c - c ^ {\prime} b), \end{array}
\]

from which the desired result follows.

\subsection{The discriminant}

DEFINITION 2. --- Let f be a monic polynomial of A{[}X{]} of degree m, and denote by E the A-algebra A{[}X{]}/(f) and by x the canonical image of X in E. We define the discriminant of f, written dis(f), as the discriminant \(D_{E/A}(1, x, \ldots, x^{m-1})\) of the basis \((1, x, \ldots, x^{m-1})\) of the A-algebra E.

For every positive integer k we write \(p_{k} = \mathrm{Tr}_{\mathrm{E}/\mathrm{A}}(x^{k})\) . By III, p.~549, Def. 2 amounts to the formula

\[
\operatorname{dis} (f) = \det \left(p _ {i + j}\right) _ {0 \leqslant i, j <   m}.\tag{45}
\]

Examples. --- 1) If \(f\) is a monic polynomial of degree 0 or 1, we have \(\operatorname{dis}(f) = 1\) , by (45).

\begin{enumerate}
\def\labelenumi{\arabic{enumi})}
\setcounter{enumi}{1}
\tightlist
\item
  Let \(f(\mathbf{X}) = \mathbf{X}^2 + \alpha \mathbf{X} + \beta\) be a monic polynomial of degree 2. Newton's relations may be written in the form (IV, p.~75)
\end{enumerate}

\[
\begin{array}{c} p _ {0} = 2 \\ p _ {1} + \alpha = 0 \\ p _ {2} + \alpha p _ {1} + 2 \beta = 0, \end{array}
\]

whence \(p_1 = -\alpha\) , \(p_2 = \alpha^2 - 2\beta\) . It follows that we have

\[
\operatorname{dis} (f) = \det \left( \begin{array}{c c} 2 & - \alpha \\ - \alpha & \alpha^ {2} - 2 \beta \end{array} \right) = \alpha^ {2} - 4 \beta .
\]

Let B be a commutative ring, \(\rho\) a homomorphism from A to B and \(\xi_{1}, \ldots, \xi_{m}\) elements of B such that

\[
{ } ^ { \rho } f ( \mathbf { X } ) = ( \mathbf { X } - \xi _ { 1 } ) \dots ( \mathbf { X } - \xi _ { m } ) .
\]

Let us write \(M\) for the matrix \((\rho(p_{i+j}))_{0 \leqslant i,j < m}\) and \(D\) for the Van der Monde matrix \((\xi_{i+1}^{j})_{0 \leqslant i,j < m}\) . By Example 4 of IV, p.~75 we have

\[
\rho (p _ {k}) = \xi_ {1} ^ {k} + \dots + \xi_ {m} ^ {k},
\]

whence \(M = {}^t D \cdot D\) ; by III, p.~532 we have \(\det D = \prod_{i > j} (\xi_i - \xi_j)\) and det \(M = (\det D)^2\) , that is

\[
\rho (\operatorname{dis} (f)) = \prod_ {i <   j} \left(\xi_ {i} - \xi_ {j}\right) ^ {2}.\tag{46}
\]

Moreover (on writing D for the derivation \(\frac{d}{d\mathbf{X}}\) ) we have

\[
D (^ {\rho} f) (\xi_ {i}) = (\xi_ {i} - \xi_ {1}) \dots (\xi_ {i} - \xi_ {i - 1}) (\xi_ {i} - \xi_ {i + 1}) \dots (\xi_ {i} - \xi_ {m})
\]

for \(1 \leqslant i \leqslant m\) , whence

\[
\rho (\operatorname{dis} (f)) = (- 1) ^ {m (m - 1) / 2} \prod_ {i \neq j} (\xi_ {i} - \xi_ {j}) = (- 1) ^ {m (m - 1) / 2} \prod_ {i = 1} ^ {m} D (^ {\rho} f) (\xi_ {i}).
\]

By Cor. 1 of IV, p.~80 applied to the universal decomposition of \(f\) , we finally obtain

\[
\operatorname{res} (f, \mathrm{D} f) = \operatorname{res} (\mathrm{D} f, f) = (- 1) ^ {m (m - 1) / 2} \operatorname{dis} (f).\tag{47}
\]

PROPOSITION 10. --- Let \(m \geqslant 1\) . There exists a unique polynomial \(\Delta \in \mathbf{Z}[\mathbf{A}_1, \ldots, \mathbf{A}_m]\) with the following property: for any commutative ring A and monic polynomial \(f = \mathbf{X}^m + \sum_{i=1}^{m} a_i \mathbf{X}^{m-i}\) in A{[}X{]} we have

\[
\operatorname{dis} (f) = \Delta \left(a _ {1}, \dots , a _ {m}\right).\tag{48}
\]

Moreover \(\Delta\) is of degree \(\leqslant 2m - 2\) and if we assign weight \(i\) to \(\mathbf{A}_i\) then \(\Delta\) is isobaric of weight \(m(m - 1)\) .

\begin{enumerate}
\def\labelenumi{\alph{enumi})}
\item
  Uniqueness of \(\Delta\) : if \(\Delta\) satisfies (48), we have in particular \(\Delta = \text{dis}(F)\) , where \(F\) is the polynomial \(X^{m} + \sum_{i=1}^{m} A_{i} X^{m-i}\) with coefficients in \(\mathbf{Z}[A_{1}, \ldots, A_{m}]\) .
\item
  Existence of \(\Delta\) : let \(s_1, \ldots, s_m\) be the elementary symmetric polynomials in the indeterminates \(\mathbf{X}_1, \ldots, \mathbf{X}_m\) . There exists a polynomial \(\Delta \in \mathbf{Z}[\mathbf{A}_1, \ldots, \mathbf{A}_m]\) , isobaric of weight \(m(m-1)\) , such that
\end{enumerate}

\[
\Delta \left(- s _ {1}, s _ {2}, \dots , (- 1) ^ {m} s _ {m}\right) = \prod_ {i <   j} \left(\mathrm{X} _ {i} - \mathrm{X} _ {j}\right) ^ {2};\tag{49}
\]

for the right-hand side is a symmetric polynomial P, homogeneous of degree \(m(m-1)\) in \(Z[X_{1},\ldots,X_{m}]\) (IV, p.~62, Th. 1). Now formula (46) means that \(\text{dis}(f)=P(f)\) , in the notation of IV, p.~74, and (48) follows directly from this.

\begin{enumerate}
\def\labelenumi{\alph{enumi})}
\setcounter{enumi}{2}
\tightlist
\item
  Degree of \(\Delta\) : the relation (47) and the definition of the resultant (IV, p.~76) imply the formula
\end{enumerate}

\[
(- 1) ^ {m (m - 1) / 2} \Delta = \det \left(a _ {i j}\right) _ {0 \leqslant i, j \leqslant 2 m - 2}\tag{50}
\]

with the following values of the \(a_{ij}\)

\[
\begin{array}{l l} a _ {0 0} = 1, \quad a _ {0, m - 1} = m, & a _ {0 j} = 0 \quad \text {if} \quad j \neq 0, \quad j \neq m - 1 \\ a _ {i j} = \mathbf {A} _ {i - j} & \text {for} \quad 1 \leqslant i \leqslant 2 m - 2, \quad 0 \leqslant j \leqslant m - 2 \\ a _ {i j} = (j - i + 1) \mathbf {A} _ {m + i - j - 1} & \text {for} \quad 1 \leqslant i \leqslant 2 m - 2, \quad m - 1 \leqslant j \leqslant 2 m - 2. \end{array}
\]

In these formulae it is understood that \(A_{0}=1\) and \(A_{i}=0\) for i\textless0 or i\textgreater m. Now (50) shows at once that \(\Delta\) is of degree \(\leqslant2m-2\) , as we had to prove.

Prop. 10 allows us to extend the definition of the discriminant to non-monic polynomials. Let \(m \geqslant 1\) be an integer, then there exists a unique homogeneous polynomial of degree \(2m - 2\) , \(\tilde{\Delta}\) say, in \(\mathbf{Z}[\mathbf{A}_0, \mathbf{A}_1, \ldots, \mathbf{A}_m]\) such that

\[
\Delta \left(\mathrm{A} _ {1}, \dots , \mathrm{A} _ {m}\right) = \tilde {\Delta} \left(1, \mathrm{A} _ {1}, \dots , \mathrm{A} _ {m}\right);\tag{51}
\]

for, since \(\Delta\) is of degree \(\leqslant 2m - 2\) , the rational fraction

\[
\mathbf {A} _ {0} ^ {2 m - 2} \Delta (\mathbf {A} _ {1} / \mathbf {A} _ {0}, \dots , \mathbf {A} _ {m} / \mathbf {A} _ {0})
\]

belongs to the subring \(\mathbf{Z}[\mathbf{A}_0, \mathbf{A}_1, \ldots, \mathbf{A}_m]\) of \(\mathbf{Q}(\mathbf{A}_0, \mathbf{A}_1, \ldots, \mathbf{A}_m)\) . If \(\mathbf{A}_i\) has weight \(i\) for \(0 \leqslant i \leqslant m\) , then \(\tilde{\Delta}\) is isobaric of weight \(m(m - 1)\) . If \(f\) is a polynomial of degree \(\leqslant m\) , say

\[
f = a _ {0} \mathrm{X} ^ {m} + a _ {1} \mathrm{X} ^ {m - 1} + \dots + a _ {m - 1} \mathrm{X} + a _ {m},
\]

we shall put

\[
\operatorname{dis} _ {m} (f) = \tilde {\Delta} \left(a _ {0}, a _ {1}, \dots , a _ {m}\right).\tag{52}
\]

When \(m = \deg f\) , we shall simply write \(\operatorname{dis}(f)\) for \(\operatorname{dis}_m(f)\) ; if \(f\) is monic, \(\operatorname{dis}(f)\) agrees with the discriminant defined in Def. 2, by (48), (51) and (52).

PROPOSITION 11. --- Let \(f\) in A{[}X{]} be of degree \(\leqslant m\) .

\begin{enumerate}
\def\labelenumi{(\roman{enumi})}
\item
  If \(\rho: \mathbf{A} \to \mathbf{B}\) is a ring homomorphism, we have \(\mathrm{dis}_m({}^{\rho}f) = \rho(\mathrm{dis}_m(f))\) .
\item
  Let \(\lambda, \alpha_{1}, \ldots, \alpha_{m}\) be elements of A. If \(f = \lambda(X - \alpha_{1}) \ldots (X - \alpha_{m})\) ; then we have
\end{enumerate}

\[
\operatorname{dis} _ {m} (f) = \lambda^ {2 m - 2} \prod_ {i <   j} (\alpha_ {i} - \alpha_ {j}) ^ {2}.\tag{53}
\]

\begin{enumerate}
\def\labelenumi{(\roman{enumi})}
\setcounter{enumi}{2}
\tightlist
\item
  Let \(a_0\) be the coefficient of \(\mathbf{X}^m\) in \(f\) ; then we have
\end{enumerate}

\[
\operatorname{res} _ {m, m - 1} (f, \mathbf {D} f) = \operatorname{res} _ {m - 1, m} (\mathbf {D} f, f) = (- 1) ^ {m (m - 1) / 2} a _ {0} \operatorname{dis} _ {m} (f).\tag{54}
\]

Assertion (i) is clear.

Since \(\tilde{\Delta}\) is homogeneous of degree \(2m - 2\) , we have

\[
\operatorname{dis} _ {m} (\lambda f) = \lambda^ {2 m - 2} \operatorname{dis} _ {m} (f)\tag{55}
\]

for every polynomial \(f \in \mathbf{A}[\mathbf{X}]\) of degree \(\leqslant m\) . Assertion (ii) now follows from formulae (46) and (55).

When f is monic of degree m, we have \(a_{0} = 1\) and Assertion (iii) is reduced to formula (47). Bearing in mind (55) and the relation

\[
\operatorname{res} _ {m, n} (\lambda f, \mu g) = \lambda^ {n} \mu^ {m} \operatorname{res} _ {m, n} (f, g),\tag{56}
\]

(IV, p.~77), we can pass from there to the case where \(a_0\) is invertible in A. Now the general case follows by Prop. 11, (i) and Lemma 5 (IV, p.~79).

COROLLARY 1. --- Let \(g \in A[X]\) and let \(n\) be a positive integer such that \(\deg g \leqslant n\) . We have

\[
\operatorname{dis} _ {m + n} (f g) = \operatorname{dis} _ {m} (f) \cdot \operatorname{dis} _ {n} (g) \cdot \operatorname{res} _ {m, n} (f, g) ^ {2}.\tag{57}
\]

By reasoning as before we are reduced to the case where f and g are monic of degrees m and n respectively. Now put \(B = E_{f} \otimes E_{g}\) , where \(E_{f}\) (resp. \(E_{g}\) ) is the universal decomposition algebra of \(f(\text{resp. } g)\) (IV, p.~73). Then A is a subring of B, and in B{[}X{]} we have the decompositions

\[
f = \prod_ {i = 1} ^ {m} \left(\mathrm{X} - \alpha_ {i}\right), \quad g = \prod_ {j = 1} ^ {n} \left(\mathrm{X} - \beta_ {j}\right).
\]

Hence we have

\[
f g = \prod_ {k = 1} ^ {m + n} \left(\mathbf {X} - \gamma_ {k}\right),
\]

with \(\gamma_{i} = \alpha_{i}\) for \(1 \leqslant i \leqslant m\) and \(\gamma_{m + j} = \beta_{j}\) for \(1 \leqslant j \leqslant n\) . We have the obvious identity

\[
\prod_ {k <   k ^ {\prime}} \left(\gamma_ {k} - \gamma_ {k ^ {\prime}}\right) = \prod_ {i <   i ^ {\prime}} \left(\alpha_ {i} - \alpha_ {i ^ {\prime}}\right) \cdot \prod_ {j <   j ^ {\prime}} \left(\beta_ {j} - \beta_ {j ^ {\prime}}\right) \cdot \prod_ {i, j} \left(\alpha_ {i} - \beta_ {j}\right).
\]

On squaring this relation we obtain (57), by (40) and (46).

COROLLARY 2. --- If \(a_{0}\) is the coefficient of \(\mathbf{X}^{m}\) in \(f\) , we have

\[
\operatorname{dis} _ {m + 1} (f) = a _ {0} ^ {2} \operatorname{dis} _ {m} (f).\tag{58}
\]

This follows from Cor. 1 on taking \(n = 1\) , \(g = 1\) , by the formula \(\operatorname{res}_{m,1}(f,1) = a_0\) (IV, p.~76, Example 1).

COROLLARY 3. --- Let A be a field, and let f be a non-constant polynomial in A{[}X{]}. For f and Df to be relatively prime, it is necessary and sufficient that \(\text{dis}(f) \neq 0\) .

This follows from Prop. 11, (iii) and Cor. 2 of IV, p.~78.

Remark. --- A double application of Cor. 2 above shows that we have \(\mathrm{dis}_{m}(f)=0\) for every polynomial f of degree \(\leqslant m-2\) .

Examples. --- 3) Let \(m = 2\) . By Example 2 (IV, p.~81) we have \(\Delta(\mathbf{A}_1, \mathbf{A}_2) = \mathbf{A}_1^2 - 4\mathbf{A}_2\) , whence \(\tilde{\Delta}(\mathbf{A}_0, \mathbf{A}_1, \mathbf{A}_2) = \mathbf{A}_1^2 - 4\mathbf{A}_0\mathbf{A}_2\) . In other words, we have

\[
\operatorname{dis} _ {2} \left(a _ {0} \mathbf {X} ^ {2} + a _ {1} \mathbf {X} + a _ {2}\right) = a _ {1} ^ {2} - 4 a _ {0} a _ {2}.
\]

\begin{enumerate}
\def\labelenumi{\arabic{enumi})}
\setcounter{enumi}{3}
\tightlist
\item
  Consider the polynomial
\end{enumerate}

\[
\mathbf {F} = \mathbf {A} _ {0} \mathbf {X} ^ {3} + 3 \mathbf {A} _ {1} \mathbf {X} ^ {2} + 3 \mathbf {A} _ {2} \mathbf {X} + \mathbf {A} _ {3}
\]

with coefficients in \(\mathbf{Q}[\mathbf{A}_0,\mathbf{A}_1,\mathbf{A}_2,\mathbf{A}_3]\) . We have

\[
\mathbf {D F} = 3 \left(\mathbf {A} _ {0} \mathbf {X} ^ {2} + 2 \mathbf {A} _ {1} \mathbf {X} + \mathbf {A} _ {2}\right),
\]

\[
\mathrm{F} - 1 / 3 \mathrm{X}. \mathrm{DF} = \mathrm{A} _ {1} \mathrm{X} ^ {2} + 2 \mathrm{A} _ {2} \mathrm{X} + \mathrm{A} _ {3}.
\]

By the formulae (54) (IV, p.~84) and (29) (IV, p.~77) we have

\[
\mathrm{A} _ {0} \cdot \operatorname{dis} _ {3} (\mathrm{F}) = - \operatorname{res} _ {2, 3} (\mathrm{DF}, \mathrm{F}) = - \operatorname{res} _ {2, 3} (\mathrm{DF}, \mathrm{F} - 1 / 3 \mathrm{X}. \mathrm{DF}).
\]

Applying Cor. 2 of IV, p.~80 we finally obtain

\[
\operatorname{dis} _ {3} (\mathrm{F}) = - 2 7 \operatorname{res} _ {2, 2} \left(\mathrm{A} _ {0} \mathrm{X} ^ {2} + 2 \mathrm{A} _ {1} \mathrm{X} + \mathrm{A} _ {2}, \mathrm{A} _ {1} \mathrm{X} ^ {2} + 2 \mathrm{A} _ {2} \mathrm{X} + \mathrm{A} _ {3}\right).
\]

By Example 3 of IV, p.80 we thus have

\[
\tilde {\Delta} \left(\mathrm{A} _ {0}, 3 \mathrm{A} _ {1}, 3 \mathrm{A} _ {2}, \mathrm{A} _ {3}\right) = - 2 7 \left(\mathrm{A} _ {0} \mathrm{A} _ {3} - \mathrm{A} _ {1} \mathrm{A} _ {2}\right) ^ {2} - 1 0 8 \left(\mathrm{A} _ {1} \mathrm{A} _ {3} - \mathrm{A} _ {2} ^ {2}\right) \left(\mathrm{A} _ {1} ^ {2} - \mathrm{A} _ {0} \mathrm{A} _ {2}\right).
\]

After some calculations we find, that if \(f = a_0\mathbf{X}^3 + a_1\mathbf{X}^2 + a_2\mathbf{X} + a_3\) , we have

\[
\operatorname{dis} _ {3} (f) = a _ {1} ^ {2} a _ {2} ^ {2} + 1 8 a _ {0} a _ {1} a _ {2} a _ {3} - 4 a _ {1} ^ {3} a _ {3} - 4 a _ {0} a _ {2} ^ {3} - 2 7 a _ {0} ^ {2} a _ {3} ^ {2}.
\]

In particular, we have

\[
\operatorname{dis} \left(\mathbf {X} ^ {3} + p \mathbf {X} + q\right) = - \left(4 p ^ {3} + 2 7 q ^ {2}\right).
\]

\BourbakiExercises

\BourbakiExerciseGroup

\begin{enumerate}
\def\labelenumi{\arabic{enumi})}
\item
  When M is a subgroup of the additive group R of real numbers, generalize the Euclidean division of polynomials in one indeterminate to the group algebra of M.
\item
  Let \(f\) be a polynomial \(\neq 0\) in \(\mathbf{A}[\mathbf{X}]\) of degree \(n\) and leading coefficient \(\alpha_0\) . Let \(\mathbf{M}\) be the submodule of \(\mathbf{A}[\mathbf{X}]\) (qua A-module) consisting of the polynomials in which the coefficient of the term of degree \(m\) (for arbitrary \(m\) in \(\mathbf{N}\) ) is divisible by \(\alpha_0^\mu\) , where \(\mu = \max(m - n + 1, 0)\) ; show that for any polynomial \(g\) in \(\mathbf{M}\) there exist two polynomials \(u, v\) in \(\mathbf{A}[\mathbf{X}]\) such that \(\deg v < n\) and \(g = uf + v\) .
\item
  \begin{enumerate}
  \def\labelenumii{\alph{enumii})}
  \tightlist
  \item
    In the algebra A{[}X{]} of polynomials in one indeterminate over A, the mapping \((u,v)\mapsto u(v)\) is an internal law of composition; show that this law is associative and distributive on the left with respect to addition and multiplication in A{[}X{]}.
  \end{enumerate}
\end{enumerate}

\begin{enumerate}
\def\labelenumi{\alph{enumi})}
\setcounter{enumi}{1}
\item
  If A is an integral domain show that the relation \(u(v)=0\) implies that u=0 or v is constant, and that when \(u\neq0\) and \(\deg v>0\) , the degree of \(u(v)\) is equal to the product of the degrees of u and v.
\item
  From now on assume that A is a field. Let u and v be two polynomials of degree \textgreater0 in A{[}X{]} and f a polynomial of degree \textgreater0; show that if q is the quotient and r the remainder in the Euclidean division of u by v, \(q(f)\) and \(r(f)\) are the quotient and remainder in the Euclidean division of \(u(f)\) by \(v(f)\) .
\item
  For each polynomial \(f\) of \(\mathbf{A}[\mathbf{X}]\) denote by \(\mathrm{I}(f)\) the set of all polynomials of the form \(u(f)\) , where \(u\) ranges over \(\mathbf{A}[\mathbf{X}]\) ; this is a subring of \(\mathbf{A}[\mathbf{X}]\) . For \(\mathrm{I}(f) = \mathrm{I}(g)\) it is necessary and sufficient that \(g = \lambda f + \mu\) , where \(\lambda \neq 0\) and \(\mu\) are in \(\mathbf{A}\) .
\item
  Let \(f\) and \(g\) be two polynomials of degree \(> 0\) in \(\mathbf{A}[\mathbf{X}]\) ; show that \(I(f) \cap I(g)\) is either equal to \(\mathbf{A}\) or of the form \(I(h)\) , where \(h\) is a polynomial of degree \(> 0\) (excluding the first possibility, consider in \(I(f) \cap I(g)\) a polynomial \(h\) of least possible degree \(> 0\) ).
\end{enumerate}

\begin{enumerate}
\def\labelenumi{\arabic{enumi})}
\setcounter{enumi}{3}
\tightlist
\item
  \begin{enumerate}
  \def\labelenumii{\alph{enumii})}
  \tightlist
  \item
    Suppose that for every non-zero \(n \in \mathbf{Z}\) , the relation \(n\xi = 0\) implies \(\xi = 0\) in \(\mathbf{A}\) . Show that the mapping \(f \mapsto f(\mathrm{D}_1, \ldots, \mathrm{D}_n)\) is an injective homomorphism of the polynomial algebra \(\mathbf{A}[\mathbf{X}_1, \ldots, \mathbf{X}_n] = \mathbf{E}\) into the algebra (over \(\mathbf{A}\) ) of endomorphisms of the \(\mathbf{A}\) -module \(\mathbf{E}\) (use induction on \(n\) ).
  \end{enumerate}
\end{enumerate}

\begin{enumerate}
\def\labelenumi{\alph{enumi})}
\setcounter{enumi}{1}
\tightlist
\item
  Let \(m\) be an integer \(\geqslant 1\) such that \(m: 1 = 0\) in \(\mathbf{A}\) . Show that we have \(\mathrm{D}_i^m = 0\) for every partial derivation \(\mathrm{D}_i\) ( \(1 \leqslant i \leqslant n\) ) in \(\mathbf{A}[\mathbf{X}_1, \ldots, \mathbf{X}_n]\) .
\end{enumerate}

\begin{enumerate}
\def\labelenumi{\arabic{enumi})}
\setcounter{enumi}{4}
\tightlist
\item
  Let \(u \in \mathbf{A}[(\mathbf{X}_i)_{i \in I}]\) and \(r\) an integer \(\geqslant 0\) , and suppose that in \(\mathbf{A}\) for each \(n > 0\) , the relation \(n\xi = 0\) implies \(\xi = 0\) . If we have the relation
\end{enumerate}

\[
\sum_ {i \in I} \left(\mathrm{D} _ {i} u\right) (\mathbf {X}) \mathbf {X} _ {i} = r u (\mathbf {X}),
\]

then \(u\) is homogeneous of degree \(r\) .

\begin{enumerate}
\def\labelenumi{\arabic{enumi})}
\setcounter{enumi}{5}
\tightlist
\item
  Let K be a commutative field, I a set and L the free associative algebra on the set I over the ring K (III, p.~448, Def. 2) with the graduation defined in III, p.~458. We define the degree of an element of L and the notion of homogeneous element as in IV, p.~2. a) Show that \(\deg(uv) = \deg(u) + \deg(v)\) and \(\deg(u + v) \leqslant \sup(\deg(u), \deg(v))\) for \(u, v \in L\) .
\end{enumerate}

\begin{enumerate}
\def\labelenumi{\alph{enumi})}
\setcounter{enumi}{1}
\item
  Given \(u, v, q, q' \in L\) such that \(\deg(u) \geqslant \deg(v)\) and \(\deg(qu - q'v) < \deg(ku)\) ; show that there exists \(q'' \in L\) such that \(\deg(u - q''v) < \deg(u)\) .
\item
  Given \(u, v, q, q' \in L\) with \(\deg u \geqslant \deg v > \deg (qu - q'v)\) ; show that there exists \(q'' \in L\) such that \(\deg (u - q''v) < \deg v\) .
\item
  Given \(u, v, u', v' \in L\) , homogeneous and such that \(uv = u'v'\) and \(\deg v \geqslant \deg v' \geqslant 0\) ; there exists \(t \in L\) such that \(u' = ut, v = tv'\) .
\item
  Given \(u, v \in L\) , homogeneous, such that \(uv = vu\) , there exists \(w \in L\) , \(a, b \in K\) and \(m, n \in K\) such that \(u = aw^m\) , \(v = bw^n\) .
\item
  Let \(u, v \in L\) be such that there exist \(a, b \in L\) , non-zero, with \(au = bv\) ; there exist \(m, d \in L\) such that the left multiples (resp. right divisors) common to \(u\) and \(v\) are left multiples (resp. right divisors) of \(m\) (resp. \(d\) ).
\end{enumerate}

\BourbakiExerciseGroup

\begin{enumerate}
\def\labelenumi{\arabic{enumi})}
\item
  Suppose A is an integral domain; let \(f\) be a polynomial of the ring \(\mathbf{A}[X_1, X_2, \ldots, X_n]\) of degree \(\leqslant k_i\) in \(X_i\) (for \(1 \leqslant i \leqslant n\) ). For each value of \(i\) ( \(1 \leqslant i \leqslant n\) ) let \(H_i\) be a set of \(k_i + 1\) elements of A. Show that if \(f(x_1, x_2, \ldots, x_n) = 0\) for each element \((x_i) \in \prod_{i=1}^{n} H_i\) , then \(f = 0\) .
\item
  Assume that A is an infinite integral domain; let \(\Phi\) be a set of polynomials \(\neq0\) in the ring \(A[X_{1},X_{2},\ldots,X_{n}]\) . Show that if the power of \(\Phi\) is strictly less than that of A, then there exists a subset H of \(A^{n}\) equipotent with A such that for each \(x=(x_{i})\in H\) and for each \(f\in\Phi\) we have \(f(x)\neq0\) .
\item
  Suppose that A is an integral domain and let B be an infinite subset of A. Show that if the polynomial \(f \in A[X]\) has degree \textgreater0, then the image of B under the polynomial function \(x \mapsto f(x)\) is equipotent with B.
\item
  Suppose that the additive group of A contains an infinite subgroup G whose non-zero elements are not divisors of zero in A. Show that the mapping \(f \mapsto \tilde{f}\) of \(A{[}X_{1}, \ldots, X_{p}{]}\) into the algebra of mappings of \(A^{p}\) into A is injective (observe that a polynomial of degree n in one indeterminate cannot have more than n distinct roots lying in G). This holds in particular when there is in A an element \(x_{0}\) , not a divisor of 0 and of infinite order in the additive group of A.
\item
  Let K be a field of infinitely many elements and let Q be the quaternion algebra over K, of type \((-1,0,-1)\) (III, p.~445). Show that the polynomial \(X^{2}+1\) has infinitely many roots in Q.
\item
  Let K be a finite field with q elements.
\end{enumerate}

\begin{enumerate}
\def\labelenumi{\alph{enumi})}
\tightlist
\item
  In the ring \(\mathbf{K}[\mathbf{X}_1, \mathbf{X}_2, \ldots, \mathbf{X}_n]\) let \(\mathfrak{a}\) be the ideal generated by the \(n\) polynomials
\end{enumerate}

\(X_{i}^{q}-X_{i}\quad(1\leqslant i\leqslant n)\) . Show that if \(f\in\mathfrak{a}\) , then \(f(x_{1},x_{2},\ldots,x_{n})=0\) for all \((x_{i})\in\mathbf{K}^{n}\) (observe that the multiplicative group \(K^{*}\) is of order q-1).

\begin{enumerate}
\def\labelenumi{\alph{enumi})}
\setcounter{enumi}{1}
\item
  If \(f\) is any polynomial of \(\mathbf{K}[\mathbf{X}_1, \mathbf{X}_2, \ldots, \mathbf{X}_n]\) , show that there exists precisely one polynomial \(\vec{f}\) such that \(\deg_i \vec{f} \leqslant q - 1\) for \(1 \leqslant i \leqslant n\) and such that \(f \equiv \vec{f} \pmod{\mathfrak{a}}\) . We thus have \(\deg \vec{f} \leqslant \deg f\) ; if \(f\) is such that \(f(x_1, x_2, \ldots, x_n) = 0\) for each \((x_i) \in \mathbf{K}^n\) , \(f\) belongs to the ideal \(\mathfrak{a}\) , which is thus the inverse image of 0 under the homomorphism \(f \mapsto \vec{f}\) (use Exercise 1).
\item
  Let \(f_1, f_2, \ldots, f_m\) be non-zero polynomials of \(\mathbf{K}[\mathbf{X}_1, \mathbf{X}_2, \ldots, \mathbf{X}_n]\) such that \(f_i(0, 0, \ldots, 0) = 0 (1 \leqslant i \leqslant m)\) and the sum of the total degrees of the \(f_i\) is \(< n\) . Show that there exists an element \((x_1, x_2, \ldots, x_n) \in \mathbf{K}^n\) distinct from \((0, 0, \ldots, 0)\) such that
\end{enumerate}

\[
f _ {i} (x _ {1}, x _ {2}, \dots , x _ {n}) = 0 \quad \text { for } \quad 1 \leqslant i \leqslant m.
\]

(Observe that if this were not so, the polynomial \(\prod_{i=1}^{m} (1 - f_i^{q-1})\) would be congruent mod a to the polynomial \(\prod_{i=1}^{n} (1 - X_i^{q-1})\) , and use \(b\) ).

\begin{enumerate}
\def\labelenumi{\arabic{enumi})}
\setcounter{enumi}{6}
\item
  Let K be a finite field with q elements, and B the product ring \(K^{1}\) where I is an infinite set. Give an example of a non-zero polynomial f of B{[}X{]} such that \(f(x)=0\) for all \(x\in B\) .
\item
  Let B be the additive group \(\mathbf{Z} \oplus (\mathbf{Z}/2\mathbf{Z})\) , with the commutative ring structure obtained by putting \((a, b) \cdot (a', b') = (aa', ab' + ba')\) for \((a, b) \in \mathbf{B}\) and \((a', b') \in \mathbf{B}\) . Let \(\varepsilon = (0, 1) \in \mathbf{B}\) , and \(f(\mathbf{X}) = \varepsilon \mathbf{X}(\mathbf{X} - 1) \in \mathbf{B}[\mathbf{X}]\) . Then \(f(\alpha) = 0\) for all \(\alpha \in B\) .
\item
  Let \(k\) be an integer \(> 0\) such that \(k \cdot 1 = 0\) in \(\mathbf{A}\) . Let \(h\) be an integer \(> 0\) and \(f\) the polynomial \((\mathbf{X} - \alpha)^{k + h} + (\mathbf{X} - \alpha)^k\) . Then \(\alpha\) is a root of order \(k\) of \(f\) and a root of order \(\geqslant k + h - 1\) of \(Df\) .
\end{enumerate}

10) Let \(f\) be a non-zero element of \(\mathbf{Z}[\mathbf{X}_1, \ldots, \mathbf{X}_n]\) . Let \(\mathrm{N}(q)\) be the number of zeros of \(f\) in the cell \(|x_1| \leqslant q, \ldots, |x_n| \leqslant q\) . Then \(\mathrm{N}(q) / q^n\) tends to 0 as \(q \to +\infty\) .

\begin{enumerate}
\def\labelenumi{\arabic{enumi})}
\setcounter{enumi}{10}
\item
  Let P be the set of all prime numbers, \(P'\) an infinite subset of P and \(f \in \mathbf{Z}[(\mathbf{X}_i)]\) . For each \(p \in P'\) let \(f_p\) be the element of \((\mathbf{Z}/p\mathbf{Z})[(\mathbf{X}_i)]\) obtained by applying to the coefficients of f the canonical homomorphism of Z onto Z/pZ. Suppose that for each \(p \in P'\) and each \(x \in (\mathbf{Z}/p\mathbf{Z})^l\) we have \(f_p(x) = 0\) ; then f = 0.
\item
  Let \(\mathbf{M}\) be an A-module, and let \(\alpha_0, \ldots, \alpha_n \in \mathbf{A}\) be such that for \(i \neq j\) the homothety with respect to \(\alpha_i - \alpha_j\) in \(\mathbf{M}\) is bijective. For any \(m_0, \ldots, m_n \in M\) there exists one and only one \(f \in \mathbf{M}[\mathbf{X}]\) such that \(f(\alpha_i) = m_i\) for \(0 \leqslant i \leqslant n\) and \(\deg f \leqslant n\) .
\item
  Let \(\alpha_{i}\) ( \(1 \leqslant i \leqslant n\) ) be \(n\) distinct elements of a commutative field \(\mathbf{K}\) , \(\beta_{i}\) ( \(1 \leqslant i \leqslant n\) ) and \(\gamma_{i}\) ( \(1 \leqslant i \leqslant n\) ) \(2n\) arbitrary elements of \(\mathbf{K}\) . Show that there exists one and only one polynomial \(f \in \mathbf{K}[\mathbf{X}]\) of degree \(\leqslant 2n - 1\) such that \(f(\alpha_{i}) = \beta_{i}\) and \(f'(\alpha_{i}) = \gamma_{i}\) for \(1 \leqslant i \leqslant n\) (Hermite interpolation formula) (to begin with consider the case where \(2n - 1\) of the \(2n\) elements \(\beta_{i}, \gamma_{i}\) are zero).
\end{enumerate}

\BourbakiExerciseGroup

\begin{enumerate}
\def\labelenumi{\arabic{enumi})}
\item
  Let E be an associative commutative and unital algebra over a commutative field K; let \(\mathbf{x} = (x_i)_{i \in I}\) be a family of pairwise permutable elements of E. Let U be the subring of \(\mathrm{K}(\mathrm{X}_i)_{i \in I}\) consisting of all \(f \in \mathrm{K}(\mathrm{X}_i)_{i \in I}\) such that x is substitutable in f.~Show that if in E the set of non-invertible elements is an ideal, then the same is true of U. Show that the ideal of non-invertible elements in U is then maximal.
\item
  \begin{enumerate}
  \def\labelenumii{\alph{enumii})}
  \tightlist
  \item
    Let K be a commutative field and \(u = a_{m}X^{m} + a_{m+1}X^{m+1} + \cdots + a_{n}X^{n}\) a polynomial of K{[}X{]} such that \(a_{m} \neq 0\) , \(a_{n} \neq 0\) ( \(0 \leq m \leq n\) ). Show that if g is a rational fraction of degree \(d \neq 0\) in K(X), then \(u(g)\) is non-zero and is of degree nd if d \textgreater{} 0, and of degree md if d \textless{} 0.
  \end{enumerate}
\end{enumerate}

\begin{enumerate}
\def\labelenumi{\alph{enumi})}
\setcounter{enumi}{1}
\tightlist
\item
  Deduce that a non-constant rational fraction g of K(X) is substitutable in every rational fraction of K(X) (observe that if g is of degree 0, there exists \(\alpha \in K\) such that \(g - \alpha\) has degree \textless{} 0).
\end{enumerate}

\begin{enumerate}
\def\labelenumi{\arabic{enumi})}
\setcounter{enumi}{2}
\tightlist
\item
  A rational fraction \(f \in K(X_1, X_2, \ldots, X_n)\) is said to be homogeneous if it is equal to the quotient of two homogeneous polynomials (the denominator being \(\neq 0\) ). Show that for \(f\) to be homogeneous it is necessary and sufficient that
\end{enumerate}

\[
f \left(\mathrm{ZX} _ {1}, \mathrm{ZX} _ {2}, \dots , \mathrm{ZX} _ {n}\right) = \mathrm{Z} ^ {d} f \left(\mathrm{X} _ {1}, \dots , \mathrm{X} _ {n}\right)
\]

where \(d\) is the degree of \(f\) .

\begin{enumerate}
\def\labelenumi{\arabic{enumi})}
\setcounter{enumi}{3}
\tightlist
\item
  Show that the Lagrange interpolation formula (IV, p.~16, Prop. 6) may be written
\end{enumerate}

\[
f (\mathbf {X}) = \omega (\mathbf {X}) \sum_ {i = 1} ^ {n} \frac {\beta_ {i}}{\omega^ {\prime} (\alpha_ {i}) (\mathbf {X} - \alpha_ {i})}
\]

where \(\omega(X) = (X - \alpha_1)(X - \alpha_2) \ldots (X - \alpha_n)\) . Deduce that if \(g\) is a polynomial of \(K[X]\) of degree \(\leqslant n - 2\) , then

\[
\sum_ {i = 1} ^ {n} \frac {g (\alpha_ {i})}{\omega^ {\prime} (\alpha_ {i})} = 0
\]

(consider the polynomial \(f(\mathbf{X}) = \mathbf{X}g(\mathbf{X})\) ).

5) Let \(\mathbf{K}\) be a commutative field of characteristic 0. Show that if a rational fraction \(u \in \mathbf{K}(\mathbf{X})\) is such that \(\mathbf{D}u = 0\) , then \(u\) is a constant.

\begin{enumerate}
\def\labelenumi{\arabic{enumi})}
\setcounter{enumi}{5}
\item
  Generalize Euler's identity to homogeneous rational fractions (Exercise 3) over a field of characteristic 0 (consider the rational fraction \(\frac{1}{Z^m} f(ZX_1, ZX_2, ..., ZX_n)\) with respect to \(Z\) and use Exercise 5).
\item
  Let \(\mathbf{K}\) be a commutative field. Show that
\end{enumerate}

\[
[ \mathbf {K} (\mathrm{T}): \mathbf {K} ] = \operatorname{Sup} (\operatorname{Card} (\mathbf {K}), \operatorname{Card} (\mathbf {N})).
\]

(Use the linear independence of the \(1 / (T - a)\) , \(a \in \mathbf{K}\) , to show that

\[
[ \mathrm{K} (\mathrm{T}): \mathrm{K} ] \geqslant \operatorname{Card} (\mathrm{K}).)
\]
